\documentclass[11pt]{article}
% BEGIN preamble.tex
\usepackage[T1]{fontenc}
\usepackage[utf8]{inputenc}
\usepackage{lmodern}
\usepackage[margin=1in,headheight=14pt]{geometry}
\usepackage{amsmath,amssymb,amsthm,mathtools,mathrsfs,bm}
\usepackage{microtype,booktabs,longtable,array,enumitem}
\usepackage{xcolor}
\definecolor{navy}{HTML}{173B57}
\definecolor{teal}{HTML}{076C78}
\usepackage[colorlinks=true,linkcolor=navy,citecolor=teal,urlcolor=teal,breaklinks=true]{hyperref}
\usepackage{bookmark}
\usepackage{fancyhdr}
\pagestyle{fancy}
\fancyhf{}
\fancyhead[L]{\small\itshape Bilinear tails and harmonic determinants}
\fancyhead[R]{\small ProveIt research continuation}
\fancyfoot[C]{\thepage}
\renewcommand{\headrulewidth}{0.3pt}
\newtheorem{theorem}{Theorem}[section]
\newtheorem{proposition}[theorem]{Proposition}
\newtheorem{lemma}[theorem]{Lemma}
\newtheorem{corollary}[theorem]{Corollary}
\newtheorem{conjecture}[theorem]{Conjecture}
\theoremstyle{definition}
\newtheorem{definition}[theorem]{Definition}
\newtheorem{example}[theorem]{Example}
\newtheorem{question}{Question}
\theoremstyle{remark}
\newtheorem{remark}[theorem]{Remark}
\numberwithin{equation}{section}
\newcommand{\dd}{\,\mathrm d}
\newcommand{\Li}{\operatorname{Li}}
\newcommand{\FP}{\operatorname{FP}}
\newcommand{\CT}{\operatorname{CT}}
\newcommand{\Res}{\operatorname{Res}}
\newcommand{\Rea}{\operatorname{Re}}
\newcommand{\Ima}{\operatorname{Im}}
\newcommand{\Real}{\operatorname{Re}}
\newcommand{\Imag}{\operatorname{Im}}
\newcommand{\Q}{\mathbb Q}
\newcommand{\R}{\mathbb R}
\newcommand{\C}{\mathbb C}
\newcommand{\Z}{\mathbb Z}
\newcommand{\N}{\mathbb N}
\newcommand{\rank}{\operatorname{rank}}
\newcommand{\ord}{\operatorname{ord}}
\newcommand{\diag}{\operatorname{diag}}
\newcommand{\adj}{\operatorname{adj}}
\newcommand{\Span}{\operatorname{span}}
\newcommand{\fall}[2]{#1^{\underline{#2}}}
\newcommand{\ee}{\boldsymbol e}
\newcommand{\code}[1]{\texttt{\detokenize{#1}}}
\setlist{itemsep=3pt,topsep=5pt}
\setlength{\emergencystretch}{2em}
\allowdisplaybreaks[2]

% END preamble.tex

\hypersetup{
 pdftitle={Bilinear Alternating Tails and Harmonic Determinants},
 pdfauthor={Research continuation prepared for Vladimir Reshetnikov},
 pdfsubject={Exact special-function identities, cutoff jets, and reconstruction}}
\begin{document}
\begin{titlepage}
\thispagestyle{empty}
\vspace*{8mm}
{\large\color{teal}\sffamily PROVEIT\quad RESEARCH CONTINUATION\par}
\vspace{14mm}
{\Huge\bfseries\color{navy}Bilinear Alternating Tails\par
\vspace{3mm}and Harmonic Determinants\par}
\vspace{8mm}
{\Large Rational certificates, Gamma identities, cutoff jets,\par
\vspace{2mm}and complete order-ten reconstruction\par}
\vspace{10mm}
{\large Prepared for Vladimir Reshetnikov\par}
\vspace{2mm}
{11 October 2026\par}
\vspace{11mm}
{\color{teal}\hrule height .8pt}
\vspace{7mm}
\noindent\textbf{Abstract.}
We develop exact identities from alternating Hurwitz tails and finite
rational moment determinants. A classical Gamma quotient gives two
bilinear endpoint laws with independent complex orders and shifts,
all mixed generalized-Stieltjes derivatives, and anchored primitives
expressed solely through Gamma ratios. We classify the rational
two-adjacent-product certificates by an integer-orbit divisor condition
and obtain explicit two-variable polygamma kernels. For the diagonal
square law we calculate every nonpositive-order cutoff contact and
every spectral jet, including the precise change under translating
the cutoff. For real rational weights, we prove that total formal
harmonic degree equals the generic coefficient-pencil rank, with only
finitely many deficient numerical cutoffs. A positive finite-difference
family realizes every rank at the same critical decay and has an exact
hierarchy of spectral poles with Gamma-product leading coefficients. A critical two-particle
cutoff has a depth-two polylogarithm generator, an explicit Euler
primitive, and an ordinary-zeta integral endpoint. Finally, an exact polynomial congruence completes all smaller-order-ten
harmonic-tail reconstructions and supplies an ordinary convergent
polygamma identity family. A cyclotomic audit identifies a previously
proved coefficient that should be removed from an open-test list.
Proofs, exact computational certificates, independent diagnostics,
source attribution, and further research targets are supplied.
\vfill
\noindent\small
\textbf{Inspected repository revision}\par
\texttt{da2792fc9cc8ef19c6a314be6c65e1437fbdc66b}\par
\vspace{3mm}
\textbf{Research status}\par
New results are distinguished from inherited inputs and open questions.
The order-ten argument includes an exact finite computer-assisted
certificate. No arithmetic-independence or exhaustive priority claim
is made.
\end{titlepage}

\tableofcontents
\clearpage
% BEGIN sections/01_scope.tex
\section{Scope, principal results, and provenance}
\label{scope:sec}

This article continues the polylogarithm and special-function programme
in ProveIt at repository revision
\texttt{da2792fc9cc8ef19c6a314be6c65e1437fbdc66b}, with particular attention
to the unified manuscript and the five research packages then present
in \path{docs/incoming}. The principal outcomes are exact identities
and structural nonvanishing theorems. Bounds enter only as tools for
one reconstruction proof.

The work has two related starting points. An alternating tail satisfies
a first-order inhomogeneous recurrence; a Gamma quotient converts it
to a homogeneous second-order recurrence with useful bilinear endpoint
identities. A finite sum of rational moments has an affine harmonic
matrix pencil; compressing onto a maximal-rank kernel determines its
highest formal harmonic degree. In each case the finite object is
established first. Infinite sums, derivatives, primitives, and Laurent
coefficients follow with the relevant limiting operation stated.

\subsection{Principal results}

\begin{enumerate}[leftmargin=*]
\item \textbf{Independent alternating tails.}
For independently varying complex orders and shifts, a classical Gamma
quotient produces both symmetric and alternating bilinear endpoint
laws (Theorem~\ref{bil:master}). Every independent mixed spectral
derivative at order one has an explicit formula in generalized
Stieltjes constants (Theorem~\ref{bil:stieltjes-all}). Anchored
integration produces a pure Gamma identity, including
\[
 \sum_{n=1}^{\infty}
 \frac{[h(n+2)-h(n+1)]^2}{h(n+1)h(n+2)}
       =3-\frac\pi2-\frac4\pi,
 \qquad h(z)=\sqrt2\,\frac{\Gamma((z+1)/2)}{\Gamma(z/2)}.
\]
\item \textbf{A complete rational certificate classification.}
Within the two-adjacent-product ansatz, the exact condition is
$P(x)/Q(x)=cS(x)S(x-1)$ with rational $S$ and outer color $\pm1$
(Theorem~\ref{rat:classification}). Its divisor criterion gives
obstructions as well as constructive certificates. A pair of adjacent
quadratics yields explicit two-variable polygamma kernels.
\item \textbf{Every negative-order cutoff jet.}
The square-law endpoint has an even inverse-power expansion. Its
coefficient $c_m(s)$ determines every contact at $s=-m$, including
all spectral derivatives, and gives an exact translation rule for a
different cutoff coordinate (Theorem~\ref{cut:contacts} and
Corollary~\ref{cut:translation}). For the natural cutoff, odd contacts
vanish; at even $m$, the derivative threshold $r>m$ is sharp.
\item \textbf{Harmonic degree equals pencil rank for real rational weights.}
The total formal degree is the generic rank over the rational cutoff
field. Only finitely many numerical cutoffs can be deficient, and a
precise kernel-compression test finds them
(Theorem~\ref{deg:main} and Corollary~\ref{deg:rational}). A signed
example shows why fixed-cutoff and generic claims differ.
\item \textbf{Every rank at the same critical decay.}
A shorter finite-difference stencil realizes every harmonic rank
$\rho=1,\ldots,r$ with decay $n^{-(2r-1)}$. Its spectral continuation
has exact poles of orders $\rho,\rho-1,\ldots,1$ at $1,2,\ldots,\rho$,
with explicit leading Laurent coefficients. The highest-order pole
has a Gamma-product coefficient
(Theorems~\ref{hier:rank} and \ref{hier:poles}). Its two-particle
cutoff has a short depth-two polylogarithm generator, an explicit
Euler primitive, and an ordinary-zeta integral endpoint
(Section~\ref{hier:gf-section}).
\item \textbf{Completion of all smaller-order-ten reconstructions.}
For every even $q>10$, the first $q$ centered even moments of the
$(10,q)$ harmonic-tail product have full formal rank $q$
(Theorem~\ref{ten:main}). A new coefficientwise Bernoulli congruence
closes the exact gap between two preceding reports. The resulting
ordinary convergent polygamma identities are given in
Corollary~\ref{ten:identity}.
\end{enumerate}

A separate audit establishes that the first complementary cyclotomic
test coefficient already follows from two proved manuscript formulas
(Proposition~\ref{audit:coefficient}). This is a correction to its
research status and cross-references, with credit retained by the
original reductions.

\subsection{What changes in the existing question list}

\begingroup\small
\begin{longtable}{@{}>{\raggedright\arraybackslash}p{.26\textwidth}>{\raggedright\arraybackslash}p{.43\textwidth}>{\raggedright\arraybackslash}p{.25\textwidth}@{}}
\toprule
Source question & Exact outcome here & Scope that remains \\
\midrule
\endfirsthead
\toprule Source question & Exact outcome here & Scope that remains \\
\midrule\endhead
Alternating-tail report, Question 1 & The order-ten exceptional
progression is eliminated; the order-eight completion is inherited
from the quartic report. & All smaller orders at most ten are now
covered; larger smaller orders remain. \\
\addlinespace
Same report, Question 4 & Necessary and sufficient rational
factorization for both adjacent products; independent quadratic
kernels. & Certificates in larger classes and with more products. \\
\addlinespace
Same report, Question 5 & Entire endpoint continuation, normally
convergent subtractions, all nonpositive-order contacts and their jets. &
Multivariable contacts beyond the diagonal square law. \\
\addlinespace
Same report, Question 6 & Independently ordered and shifted
\emph{alternating} tails, with weights supplied by the Gamma
normalization. & The question's ordinary Hurwitz product, arbitrary
outer order, and its full moving-pole catalogue remain open. \\
\addlinespace
Same report, Question 7 & Its first displayed coefficient is already
determined by manuscript identities; the implication is written out. &
The complete complementary two-resolvent generator. \\
\addlinespace
Finite-difference report, second research topic & Equality of degree
and rank for every real rational weight over $\R(N)$, with a precise
fixed-cutoff test. & Arbitrary complex weights and minimal depth
under functional relations. \\
\addlinespace
Finite-difference spectral family & All critical ranks, exact pole
orders, and highest Laurent coefficients for shorter stencils. &
General Gauss-ensemble overlapping resonances and unrestricted
higher-rank pencil classification. \\
\bottomrule
\end{longtable}
\endgroup

The source-question wording matters. In particular, a Gamma-weighted
alternating identity is substantial progress on independent parameters,
but it does not evaluate every independently weighted ordinary Hurwitz
product. Likewise the rank theorem concerns formal harmonic coordinates,
not arithmetic independence of the resulting numbers.

\subsection{Classical input and the limits of the novelty claim}

Alternating Hurwitz recurrences, their Stieltjes expansions, and the
associated Gamma quotient are classical ingredients
\cite{bil:HuKim,bil:WangHuKim}. The alternating-Hurwitz asymptotic
expansions and their spectral derivatives are already treated by
Hu--Kim \cite{HuKimAsymptotic}; the present contact theorem applies them
to a specific quadratic endpoint and computes the resulting finite-part
discrepancy. Products of zeta tails also have substantial prior
literature, including Hoffman \cite{bil:Hoffman}. Rational difference
equations and their rational solution algorithms long predate the
specialized divisor calculation used here \cite{rat:Abramov}.

The determinant work starts from the finite-difference report's rational
harmonic decomposition and its rank upper bound \cite{FiniteDifferenceIncoming}.
The order-ten proof reuses the fixed-size Bernoulli determinant and
the real-variable domination method with explicit credit
\cite{GeneralHankel,AlternatingIncoming,QuarticIncoming}. Its new
congruence, the complete coverage argument, and the resulting uniform
reconstruction statement are separated from those reused ingredients.

Thus ``new'' in this article means a proved development absent from
the inspected project corpus and not located in the targeted primary
literature search. It is not a claim that an exhaustive priority search
has excluded every previously published specialization. The proofs are
the basis for the mathematical claims. Numerical agreement is used only
to detect mistakes. No integer-relation fit is promoted to a theorem.

\subsection{Conventions and the rational harmonic decomposition}

The Hurwitz zeta function and generalized Stieltjes constants use
\[
 \zeta(1+u,a)=\frac1u+
       \sum_{j\ge0}\frac{(-1)^j\gamma_j(a)}{j!}u^j,
 \qquad \gamma_0(a)=-\psi(a).
\]
All Gamma logarithms and complex powers in right half-planes use the
holomorphic branches real on the positive axis. Shift derivatives and
spectral derivatives are distinguished. An empty determinant or product
is one. Ordinary harmonic numbers are
$H_n^{(s)}=\sum_{k=1}^{n}k^{-s}$, with $H_0^{(s)}=0$; shifted
cutoff sums are $H_N^{(q)}(a)=\sum_{n=0}^{N}(n+a)^{-q}$.

For clarity, here is the elementary decomposition behind the later
degree theorem. Given a rational weight $R$ without poles on the
nonnegative integers, expand each $x^kR(x)$, $0\le k\le2r-2$, into
its polynomial part and principal parts. Polynomial sums are rational
in $N$ by Faulhaber's formula. Group the pole shifts modulo integer
translation and choose one representative $a_\nu$ per orbit. A shifted
principal-part sum differs from $H_N^{(q)}(a_\nu)$ by finitely many
lower and upper endpoint terms, all rational in $N$. It follows exactly
that
\begin{equation}\label{scope:harmonic-decomposition}
 \left(\sum_{n=0}^{N}n^{i+j}R(n)\right)_{0\le i,j<r}
 =D_N+\sum_{\nu,q}H_N^{(q)}(a_\nu)C_{\nu,q},
\end{equation}
where $D_N$ is rational in $N$ and the symmetric matrices
$C_{\nu,q}$ do not depend on $N$. For a real rational weight,
conjugate coordinates may be replaced by real and imaginary parts.
Replacing these harmonic coordinates by independent indeterminates
defines the formal pencil and its degree. This replacement is a
bookkeeping construction; no algebraic-independence assertion is
needed or made.

For infinite expressions, ``finite part'' always has a displayed
meaning: either a Laurent coefficient in a spectral variable or the
constant term in powers and logarithms of a specified cutoff.
Different prescriptions are compared through their explicit endpoint
contacts, rather than silently identified.

% END sections/01_scope.tex

% BEGIN sections/02_bilinear_gamma.tex
% Integration-ready research section; labels bil:*; bibliography keys below.
\section{A Gamma normalization for independently shifted Hurwitz tails}
\label{bil:section}

The diagonal square law in the preceding report leaves open what happens
when both tail orders and both shifts vary independently. A normalization
by an elementary Gamma quotient gives a finite endpoint formula in all
four variables. The resulting weight is determined by the two orders;
an arbitrary additional outer power is not being evaluated here.

\subsection{The normalization and its difference equation}

For $\Re z>0$ write
\begin{align}
 \mathcal B_s(z)
 &=2^{-s}\left\{\zeta\left(s,\frac z2\right)
                    -\zeta\left(s,\frac{z+1}{2}\right)\right\},
                    \label{bil:alternating-hurwitz}\\
 \ell(z)&=\frac12\log2+\log\Gamma\left(\frac{z+1}{2}\right)
                     -\log\Gamma\left(\frac z2\right),\qquad
 h(z)=e^{\ell(z)}.\label{bil:h}
\end{align}
The difference in \eqref{bil:alternating-hurwitz} has its removable
value at $s=1$ and is entire in $s$. All logarithms of Gamma in
\eqref{bil:h} are the holomorphic branches on the right half-plane that
are real on the positive axis. Powers of $h$ mean $e^{s\ell}$.
The elementary recurrences give
\begin{equation}\label{bil:basic-recurrences}
 \mathcal B_s(z)+\mathcal B_s(z+1)=z^{-s},\qquad
 \ell(z)+\ell(z+1)=\log z,\qquad h(z)h(z+1)=z.
\end{equation}
The alternating Hurwitz recurrence and its order derivatives are
classical; see Hu--Kim \cite{bil:HuKim}. The quotient $h$ is also
classical: in the notation of Wang--Hu--Kim \cite{bil:WangHuKim},
$h(z)=\sqrt{\pi/2}/\widetilde\Gamma(z)$.

Set
\begin{equation}\label{bil:tau-v}
 \tau_s(z)=h(z)^s\mathcal B_s(z),\qquad
 V_s(z)=\frac{z^s-(z-1)^s}{h(z)^{2s}}\quad(\Re z>1).
\end{equation}
Neither $h$ nor this rewriting of $\mathcal B_s$ is asserted to be a
new special function. Their useful feature is the homogeneous equation
\begin{equation}\label{bil:homogeneous}
 \boxed{\tau_s(z-1)-\tau_s(z+1)=V_s(z)\tau_s(z).}
\end{equation}
Indeed, multiplication of the two adjacent recurrences for
$\mathcal B_s$ by the appropriate powers of $h$ gives
\[
 \tau_s(z\pm1)=h(z)^{-s}
                  -\left(\frac{h(z\pm1)}{h(z)}\right)^s\tau_s(z).
\]
Subtracting and using \eqref{bil:basic-recurrences} proves
\eqref{bil:homogeneous}.

\subsection{Two exact bilinear endpoint laws}

\begin{theorem}[Independent orders and shifts]\label{bil:master}
Let $\Re a,\Re b>0$. For $p,q\in\mathbb C$ put
\[
 u_n=\tau_p(n+a+1),\qquad v_n=\tau_q(n+b+1),\qquad
 C_n=u_nv_{n+1}+u_{n+1}v_n,\quad
 D_n=u_nv_{n+1}-u_{n+1}v_n.
\]
These definitions include $n=-1$. For every integer $N\ge0$,
\begin{align}
 &\sum_{n=0}^{N}
 \{V_p(n+a+1)+V_q(n+b+1)\}u_nv_n
       =C_{-1}-C_N,\label{bil:finite-plus}\\
 &\sum_{n=0}^{N}(-1)^n
 \{V_p(n+a+1)-V_q(n+b+1)\}u_nv_n
       =D_{-1}+(-1)^ND_N.\label{bil:finite-minus}
\end{align}
If $\Re(p+q)>0$, both series converge absolutely and locally normally,
and hence
\begin{align}
 &\sum_{n\ge0}\{V_p(n+a+1)+V_q(n+b+1)\}u_nv_n\notag\\
 &\hspace{8mm}=\tau_p(a)\tau_q(b+1)+\tau_p(a+1)\tau_q(b),
                                      \label{bil:plus}\\
 &\sum_{n\ge0}(-1)^n\{V_p(n+a+1)-V_q(n+b+1)\}u_nv_n\notag\\
 &\hspace{8mm}=\tau_p(a)\tau_q(b+1)-\tau_p(a+1)\tau_q(b).
                                      \label{bil:minus}
\end{align}
All fixed mixed derivatives in $p,q,a,b$ may be taken term by term
in this domain. The second equality also holds as an ordinary,
possibly conditional, series on $\Re(p+q)>-2$.
\end{theorem}

\begin{proof}
The homogeneous equation yields
\[
 C_{n-1}-C_n=\{V_p(n+a+1)+V_q(n+b+1)\}u_nv_n,
\]
and
\[
 D_{n-1}+D_n=\{V_p(n+a+1)-V_q(n+b+1)\}u_nv_n.
\]
The first identity telescopes. Multiplying the second by $(-1)^n$
also telescopes, since consecutive signs are opposite. This proves
both finite identities, including their lower signs.

The standard Hurwitz expansion and the Gamma quotient expansion give,
uniformly on compact sets of orders and shifts,
\begin{align}
 \mathcal B_s(z)&=\frac12z^{-s}
                       \left(1+\frac{s}{2z}+O(z^{-2})\right),\notag\\
 h(z)&=z^{1/2}\left(1-\frac1{4z}+O(z^{-2})\right),\notag\\
 \tau_s(z)&=\frac12z^{-s/2}
                       \left(1+\frac{s}{4z}+O(z^{-2})\right),
                         \label{bil:asymptotics}\\
 V_s(z)&=\frac{s}{z}+O(z^{-2}).\notag
\end{align}
These expansions hold for arbitrary bounded complex $s$, with
sufficiently many terms in the Hurwitz expansion when needed. Thus
the summands in \eqref{bil:plus} and \eqref{bil:minus} are
$O(n^{-1-\Re(p+q)/2})$, locally uniformly, and
$C_N=O(N^{-\Re(p+q)/2})$. The leading terms cancel in $D_N$, giving
$D_N=O(N^{-1-\Re(p+q)/2})$. This proves absolute normal convergence
and the endpoint limits when $\Re(p+q)>0$. Cauchy's formula gives
the derivative assertion. Finally the finite identity
\eqref{bil:finite-minus} itself proves ordinary convergence on
$\Re(p+q)>-2$, since its remaining endpoint tends to zero there.
\end{proof}

The Hurwitz and Gamma expansions invoked here are classical
\cite{bil:DLMF}; the contribution is the bilinear finite endpoint
calculus built from them. On setting $p=q=s$ and $a=b$, division
of \eqref{bil:plus} by two gives precisely
\[
 \sum_{n\ge0}\{(n+a+1)^s-(n+a)^s\}\mathcal B_s(n+a+1)^2
 =\mathcal B_s(a+1)-a^s\mathcal B_s(a+1)^2.
\]
Thus the earlier square law is a diagonal specialization of this
four-variable family.

\begin{corollary}[The first boundary contact]\label{bil:contact}
For $p+q=0$, the ordinary symmetric series in
\eqref{bil:finite-plus} converges absolutely and satisfies
\begin{align}
 &\sum_{n\ge0}\{V_p(n+a+1)+V_{-p}(n+b+1)\}
           \tau_p(n+a+1)\tau_{-p}(n+b+1)\notag\\
 &=\tau_p(a)\tau_{-p}(b+1)
        +\tau_p(a+1)\tau_{-p}(b)-\frac12.\label{bil:contact-half}
\end{align}
\end{corollary}
\begin{proof}
The $n^{-1}$ terms in the sum of the two $V$'s cancel, so the
summand is $O(n^{-2})$. Equation \eqref{bil:asymptotics} gives
$C_N\to1/2$. Now pass to the limit in \eqref{bil:finite-plus}.
\end{proof}
In particular the ordinary boundary value differs by $1/2$ from the
analytic endpoint continuation from $\Re(p+q)>0$. On the diagonal,
the division by two explains the earlier $1/4$ contact. The present
statement does not confuse those two notions of value.

\subsection{All independent Stieltjes orders}

Use the convention
\[
 \zeta(1+t,z)=\frac1t+
       \sum_{j\ge0}\frac{(-1)^j\gamma_j(z)}{j!}t^j.
\]
For $r\ge0$ define the following explicit finite expressions:
\begin{align}
 T_r(z)&=\frac{h(z)}2
  \sum_{j=0}^{r}\binom rj
  (\ell(z)-\log2)^{r-j}(-1)^j
  \left\{\gamma_j\left(\frac z2\right)
         -\gamma_j\left(\frac{z+1}{2}\right)\right\},
                                        \label{bil:T-jet}\\
 J_r(z)&=\frac1{h(z)^2}
  \left\{z(\log z-2\ell(z))^r
      -(z-1)(\log(z-1)-2\ell(z))^r\right\}.
                                        \label{bil:J-jet}
\end{align}
Here $\Re z>0$ suffices in \eqref{bil:T-jet}, and $\Re z>1$ is
used in \eqref{bil:J-jet}. Direct coefficient extraction gives
\[
 T_r(z)=\left.\partial_s^r\tau_s(z)\right|_{s=1},\qquad
 J_r(z)=\left.\partial_s^r V_s(z)\right|_{s=1}.
\]
For example $T_0(z)=h(z)\mathcal B_1(z)$ and $J_0(z)=h(z)^{-2}$.

\begin{theorem}[Every independent mixed Stieltjes jet]
\label{bil:stieltjes-all}
For $r,s\ge0$ and $\Re a,\Re b>0$, write $z_n=n+a+1$ and
$w_n=n+b+1$. Put
\begin{align*}
 X_{r,s}(z,w)&=\sum_{j=0}^{r}\binom rj
                         J_j(z)T_{r-j}(z)T_s(w),\\
 Y_{r,s}(z,w)&=\sum_{k=0}^{s}\binom sk
                         T_r(z)J_k(w)T_{s-k}(w).
\end{align*}
Then the following are ordinary absolutely convergent identities:
\begin{align}
 \sum_{n\ge0}\{X_{r,s}(z_n,w_n)+Y_{r,s}(z_n,w_n)\}
 &=T_r(a)T_s(b+1)+T_r(a+1)T_s(b),\label{bil:jets-plus}\\
 \sum_{n\ge0}(-1)^n\{X_{r,s}(z_n,w_n)-Y_{r,s}(z_n,w_n)\}
 &=T_r(a)T_s(b+1)-T_r(a+1)T_s(b).\label{bil:jets-minus}
\end{align}
\end{theorem}
\begin{proof}
Apply $\partial_p^r\partial_q^s$ to \eqref{bil:plus} and
\eqref{bil:minus} at $(p,q)=(1,1)$. The product rule gives the two
finite binomial sums exactly. Each individual scalar summand in these
sums is $O(n^{-2}\log^{r+s}(n+2))$, so no regularization or
cancellation between divergent scalar series is implicit.
\end{proof}

For the first mixed jet, $r=s=1$, the symmetric summand is explicitly
\[
 \{J_0(z)+J_0(w)\}T_1(z)T_1(w)
 +J_1(z)T_0(z)T_1(w)+J_1(w)T_1(z)T_0(w).
\]
Its sum is $T_1(a)T_1(b+1)+T_1(a+1)T_1(b)$. Thus independent
first Stieltjes coefficients at different shifts already occur in
an exact finite endpoint identity, and \eqref{bil:jets-plus}--
\eqref{bil:jets-minus} give every higher pair without an inference
from a single diagonal derivative.

\subsection{Anchored primitives and a pure Gamma square identity}

Differentiating \eqref{bil:h} shows
\begin{equation}\label{bil:h-derivative}
 \ell'(z)=\mathcal B_1(z),\qquad \tau_1(z)=h'(z).
\end{equation}
Hence the $(p,q)=(1,1)$ laws have elementary primitives in the Gamma
quotient itself. Let $A=[a_0,a_1]$ and $B=[b_0,b_1]$ be oriented
real intervals with positive endpoints, and define
\[
 U_A(t)=h(t+a_1)-h(t+a_0),\qquad
 R_A(t)=\frac1{h(t+a_0)}-\frac1{h(t+a_1)},
\]
with analogous notation for $B$. Double integration of
\eqref{bil:plus} and \eqref{bil:minus}, using
$\int h'/h^2=-1/h$, gives
\begin{align}
 \sum_{n\ge1}\{R_A(n)U_B(n)+U_A(n)R_B(n)\}
 &=U_A(0)U_B(1)+U_A(1)U_B(0),\label{bil:primitive-plus}\\
 \sum_{n\ge1}(-1)^{n-1}\{R_A(n)U_B(n)-U_A(n)R_B(n)\}
 &=U_A(0)U_B(1)-U_A(1)U_B(0).\label{bil:primitive-minus}
\end{align}
These anchored identities have no unspecified integration constants.
Uniform convergence on the compact rectangles justifies integration;
alternatively $U_A(n)=O(n^{-1/2})$ and $R_A(n)=O(n^{-3/2})$
give absolute convergence directly.

\begin{corollary}[A Gamma quotient increment square]\label{bil:gamma-square}
For positive real $a,b$,
\begin{equation}\label{bil:gamma-square-formula}
 \boxed{\sum_{n=1}^{\infty}
   \frac{\{h(n+b)-h(n+a)\}^{2}}{h(n+a)h(n+b)}
   =\{h(b)-h(a)\}\{h(b+1)-h(a+1)\}.}
\end{equation}
\end{corollary}
\begin{proof}
Take the same interval $A=B=[a,b]$ in
\eqref{bil:primitive-plus} and divide by two.
There is also a particularly short finite certificate. If
$U_n=h(n+b)-h(n+a)$, the recurrence $h(z)h(z+1)=z$ gives
\[
 U_{n-1}U_n-U_nU_{n+1}
       =\frac{U_n^2}{h(n+a)h(n+b)}.
\]
Thus the sum through $n=N$ equals $U_0U_1-U_NU_{N+1}$, and
$U_NU_{N+1}=O(N^{-1})$ tends to zero.
\end{proof}
At $(a,b)=(1,2)$ the right side is
\begin{equation}\label{bil:gamma-square-example}
 3-\frac\pi2-\frac4\pi.
\end{equation}
Indeed $h(1)=\sqrt{2/\pi}$, $h(2)=\sqrt{\pi/2}$, and
$h(3)=2\sqrt{2/\pi}$. Equation \eqref{bil:gamma-square-formula}
also extends holomorphically to $\Re a,\Re b>0$, with the algebraic
square as written and no complex conjugation.

\subsection{Scope and further targets}

The classical ingredients are the alternating Hurwitz recurrence,
Gamma shift formula, and a two-solution discrete Lagrange identity.
Products of zeta tails already have a broad literature, including
Hoffman's multiple-zeta treatment \cite{bil:Hoffman}. The claims here
are the displayed independently parameterized endpoint laws, their
explicit mixed Stieltjes coefficients, and the anchored Gamma
consequences, as additions to the inspected project material. A
literature search did not locate these particular forms; this is not
a proof of priority. No numerical independence assertion is used.

Two precise extensions remain. First, adjoining an independent outer
power to \eqref{bil:plus} generally breaks the displayed telescoper;
one should describe the additional coordinate and its moving polar
hyperplanes. Second, at $p+q=-2m$ with $m>0$, the exact upper endpoint
provides a systematic continuation scheme, but the complete
independent-order contact algebra is not established by
Corollary~\ref{bil:contact} alone.

% END sections/02_bilinear_gamma.tex

% BEGIN sections/03_rational_certificates.tex
% Integration-ready research section; labels rat:*.
\section{Exactly when a rational adjacent-tail certificate exists}
\label{rat:section}

The previous incoming report explicitly asks which pairs of polynomial
denominators admit a rational certificate made from both adjacent
tail products. The answer is a factorization criterion. It explains
why generic shifted powers require a larger coefficient class, while
also producing a simple two-variable polygamma kernel.

\subsection{The complete rational coefficient criterion}

Let $P,Q\in\mathbb C[x]\setminus\{0\}$ and write formally
\[
 E_{n-1}=E_n+\frac{(-1)^n}{P(n)},\qquad
 F_{n-1}=F_n+\frac{(-1)^n}{Q(n)}.
\]
The actual tails will later be $E_n=\sum_{k>n}(-1)^k/P(k)$ and
$F_n=\sum_{k>n}(-1)^k/Q(k)$. For a fixed nonzero number $t$ consider
\begin{equation}\label{rat:ansatz}
 K_n=t^n\{A(n)E_nF_{n+1}+B(n)E_{n+1}F_n\},\qquad
 A,B\in\mathbb C(x).
\end{equation}
By cancellation of linear-tail terms we mean cancellation of their
formal coefficients after applying the two recurrences, with $E_n,F_n$
treated as independent indeterminates. This is the precise certificate
class being classified; possible identities exploiting additional
relations between specific tails are not excluded.

\begin{theorem}[Rational certificate classification]\label{rat:classification}
A nonzero rational coefficient pair $(A,B)$ in \eqref{rat:ansatz}
eliminates both linear-tail terms if and only if
\begin{equation}\label{rat:factorization}
 t\in\{1,-1\},\qquad
 \frac{P(x)}{Q(x)}=cS(x)S(x-1)
\end{equation}
for a nonzero rational function $S$ and a nonzero constant $c$.
For a fixed admissible $t$, all coefficient pairs are proportional to
\begin{equation}\label{rat:coefficients}
 A(x)=Q(x+1)S(x),\qquad B(x)=tQ(x+1)S(x+1).
\end{equation}
If $P/Q$ is constant, the $t=-1$ certificate gives the trivial zero
identity for proportional actual tails.
\end{theorem}

\begin{proof}
Put $H=A+B$. Direct substitution into $K_n-K_{n-1}$ shows that
linear-tail cancellation is equivalent to
\begin{equation}\label{rat:coefficient-equations}
 \frac{A(x)}{Q(x+1)}=t^{-1}\frac{B(x-1)}{Q(x)},\qquad
 \frac{B(x)}{P(x+1)}=t^{-1}\frac{A(x-1)}{P(x)}.
\end{equation}
Neither coefficient can vanish unless both do. Eliminating $B$ yields
\[
 \frac{A(x)}{A(x-2)}=
 t^{-2}\frac{Q(x+1)P(x)}{Q(x)P(x-1)}.
\]
Both rational ratios apart from $t^{-2}$ tend to one at infinity,
so $t^2=1$. Put $R=P/Q$ and $S(x)=A(x)/Q(x+1)$. The remaining
equation is
\[
 \frac{S(x)}{S(x-2)}=\frac{R(x)}{R(x-1)}.
\]
Consequently $R(x)/(S(x)S(x-1))$ is a rational function of period
one. A periodic rational function is constant: a pole would generate
infinitely many poles, and a nonconstant polynomial cannot be
periodic. This proves the necessity of \eqref{rat:factorization}.
Conversely that factorization and \eqref{rat:coefficients} satisfy
\eqref{rat:coefficient-equations} by substitution. The quotient of
two solutions for $A$ is rational and of period two, so it is
constant, proving uniqueness up to scalar.
\end{proof}

This elementary argument is a specialized rational difference-equation
calculation, within the classical subject of rational summation and
recurrence solutions \cite{rat:Abramov}. It does not claim a new general
algorithm for rational recurrences.

\begin{corollary}[Root-orbit criterion]\label{rat:root-orbits}
On each orbit $\alpha+\mathbb Z$ let $m_j$ be the multiplicity of
$x=\alpha+j$ as a zero of $P/Q$, with negative multiplicities for
poles. A rational $S$ in \eqref{rat:factorization} exists if and only if
\begin{equation}\label{rat:root-balance}
 \sum_{j\in\mathbb Z}(-1)^jm_j=0
 \quad\hbox{on every such orbit}.
\end{equation}
\end{corollary}
\begin{proof}
If $s_j$ are the multiplicities for $S$, factorization is equivalent
to $m_j=s_j+s_{j-1}$. Starting below the finite support of $(m_j)$,
the unique solution is
$s_j=\sum_{k\le j}(-1)^{j-k}m_k$. It has finite support exactly
under \eqref{rat:root-balance}. The remaining nonzero multiplicative
constant is absorbed into $c$.
\end{proof}

For example, for positive integer orders $p,q$ and shifted powers
$P(x)=(x+a)^p$, $Q(x)=(x+b)^q$, the criterion holds exactly when
\begin{equation}\label{rat:powers-obstruction}
 p=q\quad\hbox{and}\quad a-b\in2\mathbb Z.
\end{equation}
If $a-b$ is not an integer the two roots occupy separate orbits.
If it is an integer $d$, the condition is $p=(-1)^dq$, yielding
\eqref{rat:powers-obstruction}. This is an obstruction only to
\eqref{rat:ansatz}, not to special-function endpoint evaluations in
larger coefficient classes. The Gamma normalization in
Section~\ref{bil:section} supplies just such a larger class.

\subsection{Convergence and the finite endpoint}

Assume $\deg P,\deg Q\ge1$, no positive integer zeros, and an
admissible factorization. Put
\[
 H_t(x)=Q(x+1)\{S(x)+tS(x+1)\},\qquad
 W_t(n)=t^n\{H_t(n)-tH_t(n-1)\}.
\]
The certificate calculation gives exactly
\begin{equation}\label{rat:finite-difference}
 K_n-K_{n-1}=W_t(n)E_nF_n.
\end{equation}
Suppose also that the displayed coefficient values are finite on the
summation lattice and at the lower endpoint; removable values may be
taken. To state the endpoint without alternating signs in the tails,
write
\[
 e_n=\sum_{j\ge0}\frac{(-1)^j}{P(n+1+j)},\qquad
 f_n=\sum_{j\ge0}\frac{(-1)^j}{Q(n+1+j)}.
\]
Then, whenever $S(0)$ is finite and nonzero,
\begin{align}
 \sum_{n\ge0}W_t(n)e_nf_n
 &=S(0)e_0+\frac{t}{cS(0)}f_0
       -tH_t(-1)e_0f_0.\label{rat:endpoint}
\end{align}
The finite version subtracts
\[
 t^N\{A(N)e_Nf_{N+1}+B(N)e_{N+1}f_N\}
\]
from the right side of \eqref{rat:endpoint}. This also gives a
safe lower-cutoff version if one of the coefficient functions has a
pole at one of the first lattice points.

For completeness, write $p=\deg P$, $q=\deg Q$. Factorization forces
$d=(p+q)/2$ to be an integer at least one. Alternating summation
gives $e_n=O(n^{-p})$, $f_n=O(n^{-q})$, while $A,B=O(n^d)$.
For $t=1$, the difference of $H_1$ is $O(n^{d-1})$. For $t=-1$,
$S(n)-S(n+1)$ lowers the degree by at least one, so $H_{-1}$ and
its adjacent sum are $O(n^{d-1})$. Thus every series in
\eqref{rat:endpoint} is absolutely convergent, and its upper endpoint
is $O(N^{-d})$. For positive-order rational tails all endpoints
reduce by partial fractions to finite polygamma values at half-shifts.

\subsection{Two independently shifted adjacent quadratic tails}

Define, for positive integer $p$,
\begin{equation}\label{rat:D-definition}
 D_p(z)=\mathcal B_p(z)-\mathcal B_p(z+1),\qquad
 D(z)=D_1(z)=2\mathcal B_1(z)-\frac1z.
\end{equation}
In particular
\[
 D(n+x)=\sum_{j\ge0}\frac{(-1)^j}
                          {(n+x+j)(n+x+j+1)}.
\]
The denominators $P(k)=(k+x-1)(k+x)$ and
$Q(k)=(k+y-1)(k+y)$ satisfy \eqref{rat:factorization} with
$S(k)=(k+x)/(k+y)$ and $c=1$. This works for independent $x,y$;
the roots within each quadratic are adjacent.

\begin{theorem}[Finite two-variable polygamma kernels]\label{rat:kernels}
For $\Re x,\Re y>0$, both series below converge absolutely and
locally normally:
\begin{align}
 \sum_{n\ge0}(2n+x+y)D(n+x)D(n+y)
 &=\frac12\left\{\frac{x}{y}D(x)+\frac{y}{x}D(y)
                -(2xy-x-y)D(x)D(y)\right\},\label{rat:ordinary}\\
 \sum_{n\ge0}(-1)^nD(n+x)D(n+y)
 &=\frac12D(x)D(y)
    +\frac{x^2D(x)-y^2D(y)}{2xy(x-y)}.\label{rat:alternating}
\end{align}
The quotient at $x=y$ is interpreted by continuity. In particular,
\begin{equation}\label{rat:diagonal}
 \sum_{n\ge0}(-1)^nD(n+x)^2
     =\frac12D(x)^2+\frac{D(x)}x+\frac12D'(x).
\end{equation}
\end{theorem}

\begin{proof}
For the stated $P,Q,S$, the two rational coefficients are
\[
 A(n)=(n+x)(n+y+1),\qquad
 B(n)=t(n+y)(n+x+1).
\]
For $t=1$, $H_1(n)-H_1(n-1)=2(2n+x+y)$ and
$H_1(-1)=2xy-x-y$. Substitution in \eqref{rat:endpoint} and
division by two proves \eqref{rat:ordinary}. For $t=-1$,
$H_{-1}(n)=x-y$ and $W_{-1}(n)=2(x-y)(-1)^n$.
The resulting endpoint is
\[
 \frac{x}{y}D(x)-\frac{y}{x}D(y)+(x-y)D(x)D(y),
\]
which proves \eqref{rat:alternating} for $x\ne y$.
The locally normal series has a limit on the diagonal. Taking that
limit on the right gives \eqref{rat:diagonal}. The estimates
$D(n+x)=O(n^{-2})$ prove all claimed convergence statements.
\end{proof}

For example, with $L=\log2$,
\begin{align}
 \sum_{n\ge0}(-1)^nD(n+1)^2&=2L^2-\frac{\pi^2}{12},\label{rat:example1}\\
 \sum_{n\ge0}(-1)^nD(n+1)D(n+2)&=1-2L^2,\label{rat:example2}\\
 \sum_{n\ge0}(2n+3)D(n+1)D(n+2)&=2(1-L)^2.\label{rat:example3}
\end{align}
These follow from $D(1)=2L-1$, $D(2)=3/2-2L$, and
$D'(1)=1-\pi^2/6$.

\subsection{Every mixed integer order at independent shifts}

Let $K(x,y)$ and $M(x,y)$ denote the right sides of
\eqref{rat:alternating} and \eqref{rat:ordinary}, respectively.
For $p,q\ge1$ set
\[
 \mathfrak D_{p,q}=
 \frac{(-1)^{p+q-2}}{(p-1)!(q-1)!}
       \partial_x^{p-1}\partial_y^{q-1}.
\]
The shift derivative formula gives
$D^{(r)}(z)=(-1)^rr!D_{r+1}(z)$. Therefore
\begin{equation}\label{rat:all-orders}
 \boxed{\sum_{n\ge0}(-1)^nD_p(n+x)D_q(n+y)
                    =\mathfrak D_{p,q}K(x,y).}
\end{equation}
The right side is a finite expression in polygamma values and a
confluent divided difference; it contains no tail sum. Likewise
\begin{align}
 &\sum_{n\ge0}\bigl\{(2n+x+y)D_p(n+x)D_q(n+y)\notag\\
 &\qquad-\boldsymbol1_{p\ge2}D_{p-1}(n+x)D_q(n+y)
          -\boldsymbol1_{q\ge2}D_p(n+x)D_{q-1}(n+y)\bigr\}
            =\mathfrak D_{p,q}M(x,y).\label{rat:all-orders-ordinary}
\end{align}
When an indicator is zero its corresponding term is absent. The
formula follows by differentiating the linear weight as well as both
tails, so no lower-order correction has been omitted.

The identities hold meromorphically away from the nonpositive-integer
coordinate poles, with the diagonal $x=y$ removable. Their proof uses
only rational cancellation and differentiation of normally convergent
series. The obstruction \eqref{rat:powers-obstruction} and the positive
kernel family \eqref{rat:all-orders} concern different denominator
classes and are fully compatible.

% END sections/03_rational_certificates.tex

% BEGIN sections/04_cutoff_jets.tex
\section{All cutoff jets of the alternating square law}
\label{cut:sec}

The preceding independent-order identities grow from a finite endpoint
identity. This section resolves Question 5 of the incoming alternating-tail
report \cite{AlternatingIncoming}: continue its ordinary square law through
all nonpositive orders, calculate every cutoff contact, and distinguish the
continued value from the constant term of the original finite sums. The
alternating Hurwitz asymptotic expansion itself is classical; in particular,
Hu and Kim give its remainder theory and spectral derivatives
\cite{HuKimAsymptotic}. The contribution here is the explicit quadratic
endpoint polynomial, its all-index contact formula, its sharp derivative
threshold, and the resulting Gamma identities.

\subsection{A finite identity valid at every spectral order}

Throughout this section $\Rea a>0$. All powers and logarithms use the
holomorphic branch in the right half-plane. Recall
\begin{equation}\label{cut:B}
 \mathcal B_s(x)=2^{-s}\left\{\zeta\left(s,\frac{x}{2}\right)
                    -\zeta\left(s,\frac{x+1}{2}\right)\right\}.
\end{equation}
It is entire in $s$, with the common Hurwitz pole at one cancelled.
Put
\begin{align}
 F_s(x)&=x^s\mathcal B_s(x)\mathcal B_s(x+1),\label{cut:F}\\
 g_s(x)&=\{x^s-(x-1)^s\}\mathcal B_s(x)^2,\label{cut:g}\\
 S_N(s,a)&=\sum_{n=0}^{N}g_s(n+a+1),\qquad X=N+a+1.
 \label{cut:SN}
\end{align}

\begin{lemma}[Exact endpoint identity]\label{cut:finite}
For every complex $s$ and every integer $N\ge0$,
\begin{equation}\label{cut:finiteeq}
 S_N(s,a)=F_s(a)-F_s(X).
\end{equation}
For $\Rea s>0$, the original series converges absolutely and locally
normally to $F_s(a)$. Its unique entire continuation as a function of $s$
is therefore $F_s(a)$.
\end{lemma}
\begin{proof}
The recurrence
$\mathcal B_s(x)+\mathcal B_s(x+1)=x^{-s}$ gives
\[
 F_s(x-1)-F_s(x)
 =\{x^s-(x-1)^s\}\mathcal B_s(x)^2.
\]
Telescoping proves the finite formula without any limiting argument.
For $\Rea s>0$, the standard tail expansion gives
$F_s(X)=O(X^{-\Rea s})$ and
$g_s(n+a+1)=O(n^{-\Rea s-1})$, uniformly on compact parameter sets.
This proves the convergence assertion. Both sides of the finite identity
are entire in $s$, so it already applies at every negative integer.
\end{proof}

An entire continuation is a function; a cutoff finite part is a summation
prescription. Their agreement on a convergence half-plane does not make
them identical at every exceptional parameter. We calculate the discrepancy
instead of adopting an implicit convention.

\subsection{The endpoint has only even inverse powers}

Let $E_j(x)$ be the Euler polynomials, normalized by
\begin{equation}\label{cut:Euler}
 \frac{2e^{xt}}{e^t+1}=\sum_{j\ge0}E_j(x)\frac{t^j}{j!}.
\end{equation}
Thus $E_0(0)=1$, $E_1(0)=-1/2$, and $E_{2j}(0)=0$ for $j\ge1$.
For $j\ge0$ define the rational polynomials
\begin{equation}\label{cut:a}
 a_j(s)=\frac{(-1)^jE_j(0)}{2j!}(s)_j,
 \qquad (s)_0=1.
\end{equation}
Define $c_0(s)=1/4$ and, for $k\ge1$,
\begin{equation}\label{cut:c}
 c_k(s)=-\sum_{j=1}^{k-1}a_j(s)a_{k-j}(s).
\end{equation}
In particular $c_{2j+1}=0$. Every nonzero $c_{2j}$ has degree $2j$.
The first three nonconstant polynomials are
\begin{align}
 c_2(s)&=-\frac{s^2}{16},\label{cut:firstc}\\
 c_4(s)&=\frac{s^2(s+1)(s+2)}{96},\notag\\
 c_6(s)&=-\frac{s^2(s+1)(s+2)(17s^2+99s+154)}{11520}.\notag
\end{align}

\begin{proposition}[Even endpoint expansion]\label{cut:expansion}
For every $K\ge0$,
\begin{equation}\label{cut:expansioneq}
 F_s(x)=\sum_{j=0}^{K}c_{2j}(s)x^{-s-2j}
           +O\bigl(x^{-\Rea s-2K-2}\bigr)
\end{equation}
as $x\to+\infty$, locally uniformly in $s$ and in bounded right-half-plane
shifts of $x$. Every fixed spectral derivative has the differentiated
expansion, with the corresponding finite powers of $\log x$. A derivative
in $x$ improves the remainder by one inverse power.
\end{proposition}
\begin{proof}
The alternating Hurwitz expansion is
\[
 \mathcal B_s(x)\sim x^{-s}\sum_{j\ge0}a_j(s)x^{-j}.
\]
One may obtain the needed uniform remainder by applying the ordinary
Hurwitz Euler--Maclaurin expansion to the two terms in \eqref{cut:B},
using more terms than required and cancelling their common pole at one.
The usual remainder integral is uniform on each compact $s$-set after
enough terms have been retained \cite{DLMFZeta,HuKimAsymptotic}.
Its formal coefficients are also determined uniquely by
$\mathcal B_s(x)+\mathcal B_s(x+1)=x^{-s}$: the formal operator
$(1+e^{\partial_x})^{-1}$ acting on $x^{-s}$ gives \eqref{cut:a}.
Cauchy's formula on a slightly larger compact set justifies every fixed
spectral derivative of the remainder. The derivative identity
$\partial_x\mathcal B_s(x)=-s\mathcal B_{s+1}(x)$ supplies the
corresponding derivative remainder in $x$.

Write $A=x^s\mathcal B_s(x)$. By the same recurrence,
\[
 F_s(x)=x^{-s}(A-A^2)
       =x^{-s}\left\{\frac14-\left(A-\frac12\right)^2\right\}.
\]
All positive even coefficients of $A$ vanish. Thus $A-1/2$ has only
odd inverse powers, its square has only even inverse powers, and
\eqref{cut:c} gives every coefficient. Multiplying finite expansions with
their uniform remainders proves the assertion, rather than assuming
convergence of the infinite asymptotic series.
\end{proof}

\subsection{A convergent subtraction on every half-plane}

For $K\ge0$ set
\begin{equation}\label{cut:PK}
 P_K(s,x)=\sum_{j=0}^{K}c_{2j}(s)x^{-s-2j}.
\end{equation}

\begin{theorem}[Normally convergent continuation]\label{cut:normal}
On $\Rea s>-2K-2$, $\Rea a>0$, the following series converges absolutely
and locally normally, including after every fixed parameter derivative:
\begin{align}
 &\sum_{n=0}^{\infty}\bigl\{
    g_s(n+a+1)-P_K(s,n+a)+P_K(s,n+a+1)\bigr\}\notag\\
 &\hspace{30mm}=F_s(a)-P_K(s,a).
 \label{cut:normaleq}
\end{align}
Consequently adding $P_K(s,a)$ gives the same entire continuation
$F_s(a)$ on all these overlapping half-planes.
\end{theorem}
\begin{proof}
Subtract the finite telescoping difference of $P_K$ from
\eqref{cut:finiteeq}. The remaining upper endpoint is
$F_s(X)-P_K(s,X)=O(X^{-\Rea s-2K-2})$ and tends to zero on the stated
half-plane. Each summand is a unit difference of this remainder; the
$x$-derivative version of Proposition~\ref{cut:expansion} bounds it by
$O(n^{-\Rea s-2K-3})$. Uniformity on compact sets and Cauchy's formula
give local normal convergence and all fixed parameter derivatives.
\end{proof}

This is an ordinary convergent-series realization of the continued
function. It explicitly specifies what was subtracted. It neither assigns
an ordinary value to a divergent unmodified series nor discards the exact
upper endpoint.

\subsection{Every nonpositive-order contact and every spectral derivative}

For a fixed nonnegative integer $m$ and derivative order $r$, expand
$\partial_s^rS_N(s,a)|_{s=-m}$ in powers of $X=N+a+1$ and $\log X$.
Its \emph{natural-cutoff finite part} is the coefficient of
$X^0(\log X)^0$. Denote it by $\FP_X$.
The preceding expansion supplies such an asymptotic expansion to any
fixed order, with a remainder tending to zero after enough terms.

\begin{theorem}[Complete contact formula]\label{cut:contacts}
For every $m,r\ge0$ and $\Rea a>0$,
\begin{equation}\label{cut:contacteq}
 \boxed{\displaystyle
 \FP_X\left.\partial_s^rS_N(s,a)\right|_{s=-m}
   =\left.\partial_s^rF_s(a)\right|_{s=-m}-c_m^{(r)}(-m).}
\end{equation}
The contacts have the following exact properties:
\begin{enumerate}
\item If $m$ is odd, every natural-cutoff contact is zero.
\item If $m$ is even, the contact vanishes for $r>m$, and the threshold
is sharp:
\begin{equation}\label{cut:top}
 c_m^{(m)}(-m)=-\frac12 E_{m+1}(0)\ne0.
\end{equation}
\item At derivative order zero,
\begin{equation}\label{cut:Eulercontact}
 c_m(-m)=\frac14[x^m]\{E_m(x)E_m(x+1)\}.
\end{equation}
For even $m=2j\ge2$ this becomes
\begin{equation}\label{cut:Eulerconvolution}
 c_{2j}(-2j)
 =-\frac14\sum_{\substack{1\le k<2j\\k\text{ odd}}}
     \binom{2j}{k}^{\!2}E_k(0)E_{2j-k}(0),
\end{equation}
which is nonzero and has sign $(-1)^j$.
\end{enumerate}
\end{theorem}
\begin{proof}
Differentiate the endpoint expansion before setting $s=-m$.
For each $k$,
\[
 \left.\partial_s^r\{c_k(s)X^{-s-k}\}\right|_{s=-m}
 =X^{m-k}\sum_{\ell=0}^{r}\binom r\ell
       c_k^{(\ell)}(-m)(-\log X)^{r-\ell}.
\]
Only $k=m$ has exponent zero, and only $\ell=r$ has logarithmic degree
zero. The finite identity therefore gives \eqref{cut:contacteq}.
If $m$ is odd, $c_m$ is identically zero. Degree at most $m$ gives the
vanishing for $r>m$.

For the exact top coefficient, replace each $(s)_j$ in \eqref{cut:a}
by its leading monomial $s^j$. The resulting generating series is
\[
 A(t)=\frac1{1+e^{-t}}=\frac12+\frac12\tanh(t/2).
\]
The endpoint's leading-coefficient series is
\[
 A(t)-A(t)^2=A'(t)=\frac1{4\cosh^2(t/2)}.
\]
Its coefficient of $t^m$ is
$(-1)^{m+1}E_{m+1}(0)/(2m!)$; for even $m$ this proves
\eqref{cut:top}. Odd Euler values at zero are nonzero, for example by
their elementary Bernoulli-number formula.

Finally $\mathcal B_{-m}(x)=E_m(x)/2$ is the classical Euler
specialization \cite{HuKimAsymptotic}. Hence
$F_{-m}(x)=x^{-m}E_m(x)E_m(x+1)/4$.
Its constant coefficient gives \eqref{cut:Eulercontact}. For even
$m\ge2$, use $E_m(x+1)=2x^m-E_m(x)$ and $E_m(0)=0$ to obtain
\eqref{cut:Eulerconvolution}. Each summand in that convolution has
sign $(-1)^{j+1}$, so its negative quarter has the asserted nonzero sign.
\end{proof}

\begin{center}
\begin{tabular}{@{}rrrrrr@{}}
\toprule
$m$&0&2&4&6&8\\
\midrule
$c_m(-m)$&$1/4$&$-1/4$&$1$&$-43/4$&$230$\\
$c_m^{(m)}(-m)$&$1/4$&$-1/8$&$1/4$&$-17/16$&$31/4$\\
\bottomrule
\end{tabular}
\end{center}

At $s=0$ every unmodified summand is exactly zero, whereas
$F_0(a)=1/4$. Formula~\eqref{cut:contacteq} gives $0=1/4-1/4$.
At $s=-2$, the summand simplifies to $-(2n+2a+1)/4$, so
\[
 S_N(-2,a)=\frac{a^2-X^2}{4},\qquad
 F_{-2}(a)=\frac{a^2-1}{4}.
\]
The natural-cutoff constant is $a^2/4$, exactly the continued value
minus the contact $-1/4$. The first two derivative contacts at this
same order are $c_2'(-2)=1/4$ and $c_2''(-2)=-1/8$.

\subsection{The cutoff coordinate is part of the formula}

The usual integer cutoff $N$ and the natural cutoff $X=N+a+1$ have
different constant terms. The conversion is finite and explicit.
Let $T=X-h$ for a fixed complex number $h$. Define the degree-at-most-$m$
polynomial in $s$
\begin{equation}\label{cut:translated}
 C_m(s;h)=\sum_{k=0}^{m}c_k(s)
                 \binom{-s-k}{m-k}h^{m-k}.
\end{equation}
The generalized binomial here is its polynomial expression, so its
spectral derivatives at integers are unambiguous.

\begin{corollary}[Translated-cutoff contacts]\label{cut:translation}
For all $m,r\ge0$,
\begin{equation}\label{cut:translationeq}
 \FP_T\left.\partial_s^rS_N(s,a)\right|_{s=-m}
 =\left.\partial_s^rF_s(a)\right|_{s=-m}
  -\left.\partial_s^r C_m(s;h)\right|_{s=-m}.
\end{equation}
Take $h=a+1$ to obtain the finite part in powers of $N$.
Every translated contact still vanishes for $r>m$.
At $r=0$ its value is
$\sum_{k=0}^{m}c_k(-m)h^{m-k}$.
\end{corollary}
\begin{proof}
Expand $X^{-s-k}=T^{-s-k}(1+h/T)^{-s-k}$ by its finite asymptotic
binomial expansion. At $s=-m$, the unique term capable of contributing
a constant has binomial index $m-k$. Differentiating first retains its
spectral coefficient and logarithms; the logarithm-free constant is
exactly the derivative in \eqref{cut:translationeq}. Terms with $k>m$
already decay. The polynomial degree bound follows from
$\deg c_k\le k$ and
$\deg_s\binom{-s-k}{m-k}=m-k$.
\end{proof}

For example the natural contact at $s=-1$ is zero but the contact in
$T=X-h$ is $h/4$. Thus ``odd negative orders have no contact'' is
correct only with the displayed natural coordinate. Translation does
not preserve that parity assertion.

One can also write the highest derivative in closed form. Define
Euler polynomials of order two by
\[
 \frac{4e^{xt}}{(e^t+1)^2}
 =\sum_{m\ge0}E_m^{(2)}(x)\frac{t^m}{m!}.
\]
The same leading-coefficient calculation gives
\begin{equation}\label{cut:translatedtop}
 \partial_s^m C_m(s;h)=\frac{(-1)^m}{4}E_m^{(2)}(h+1).
\end{equation}
This includes \eqref{cut:top} when $h=0$ and $m$ is even.

\subsection{Explicit Gamma identities at order zero}

The classical Lerch formula gives
\begin{equation}\label{cut:Lerch}
 \left.\partial_s\mathcal B_s(x)\right|_{s=0}
 =-\log h(x),\qquad
 h(x)=\sqrt2\,\frac{\Gamma((x+1)/2)}{\Gamma(x/2)}.
\end{equation}
Put
\begin{equation}\label{cut:rho}
 \rho(x)=\frac12\log x-\log h(x)
 =\log\left\{\sqrt{\frac{x}{2}}\,
          \frac{\Gamma(x/2)}{\Gamma((x+1)/2)}\right\}.
\end{equation}
The normalized endpoint immediately gives
\begin{align}
 \left.\partial_sF_s(a)\right|_{s=0}
    &=-\frac14\log a,\label{cut:Fzero}\\
 \left.\partial_s^2F_s(a)\right|_{s=0}
    &=\frac14\log^2a-2\rho(a)^2.\notag
\end{align}
Indeed, with $D_s(a)=a^s\mathcal B_s(a)-1/2$,
$F_s(a)=a^{-s}(1/4-D_s(a)^2)$, $D_0=0$, and $D'_0=\rho(a)$.
In particular the natural finite part of the second spectral derivative
at $a=1,s=0$ is $-\tfrac12\log^2(\pi/2)$.

There is an ordinary convergent version with a useful polarization.
For $\Rea a,\Rea b>0$ write
\[
 \delta_a(n)=\log\frac{n+a+1}{n+a},\qquad
 \rho_a(n)=\rho(n+a).
\]

\begin{corollary}[A bilinear logarithmic Gamma sum]\label{cut:Gamma}
For $\Rea a,\Rea b>0$,
\begin{align}
 \sum_{n=0}^{\infty}\bigl\{
   \delta_a(n)\rho_b(n)+\delta_b(n)\rho_a(n)
        -\tfrac12\delta_a(n)\delta_b(n)\bigr\}
 =2\rho(a)\rho(b).
 \label{cut:Gammaeq}
\end{align}
The convergence is absolute and locally normal. At $a=b$ this reduces to
\begin{equation}\label{cut:Gammasquare}
 \sum_{n=0}^{\infty}\delta_a(n)
       \left\{\rho_a(n)-\frac14\delta_a(n)\right\}=\rho(a)^2.
\end{equation}
In particular the value at $a=1$ is $\tfrac14\log^2(\pi/2)$.
\end{corollary}
\begin{proof}
The recurrence $h(x)h(x+1)=x$ implies
$\rho(x)+\rho(x+1)=\tfrac12\log((x+1)/x)$.
Multiply the two shifted recurrences and rearrange:
\[
 2\{\rho_a(n)\rho_b(n)-\rho_a(n+1)\rho_b(n+1)\}
 =\delta_a(n)\rho_b(n)+\delta_b(n)\rho_a(n)
            -\tfrac12\delta_a(n)\delta_b(n).
\]
Telescoping proves the formula. Stirling's expansion gives
$\rho(n+a)=O(n^{-1})$ with a unit difference $O(n^{-2})$, so the
summand is $O(n^{-3})$ locally uniformly and the upper endpoint vanishes.
The square case follows by division by two. Equivalently it is the
second derivative of the normally subtracted identity
\eqref{cut:normaleq} at $s=0$.
\end{proof}

Every fixed shift derivative of \eqref{cut:Gammaeq} is an ordinary
convergent identity with rational functions, logarithms and polygamma
functions. For higher spectral derivatives at zero or at $-m$, the
finite formula \eqref{cut:contacteq} uses only finitely many Hurwitz
order derivatives at the same integer. No independent arithmetic
reduction of those order derivatives is asserted.

% END sections/04_cutoff_jets.tex

% BEGIN sections/05_degree_rank.tex
\section{The harmonic degree is generically equal to the pencil rank}
\label{deg:sec}

The finite-difference determinant report proves an upper bound for the
harmonic degree and asks when that bound is strict.  For real rational
weights, the answer is particularly clean: it is never strict over the
field of rational functions in the cutoff.  At a specified integer
cutoff it can be strict, but only at finitely many cutoffs.  Positivity
removes these exceptions once enough nodes have been included.

\subsection{An affine symmetric-pencil lemma}

Let $C_1,\ldots,C_t$ be real symmetric $r\times r$ matrices, write
\[
 C(X)=\sum_{j=1}^tX_jC_j,
 \qquad \rho=\operatorname{rank}_{\mathbb R(X)} C(X),
 \qquad q=r-\rho,
\]
and let $D$ be another real symmetric matrix.  The variables $X_j$ are
formal independent variables.  Their introduction makes no assertion
about algebraic independence of special functions.

\begin{lemma}[Compression and the highest degree]
\label{deg:compression}
The polynomial $P(X)=\det(D+C(X))$ has total degree at most $\rho$.
Choose any $v\in\mathbb R^t$ such that $C(v)$ has rank $\rho$, and let
$K$ be a real $r\times q$ matrix whose columns form its kernel.
Then
\begin{equation}\label{deg:isotropic}
 K^{\mathsf T}C_jK=0\qquad(1\le j\le t).
\end{equation}
If $K^{\mathsf T}DK$ is nonsingular, $P$ has total degree exactly $\rho$.
For a generic maximal-rank direction, nonsingularity of this compression
is also necessary.
\end{lemma}
\begin{proof}
In each determinant term, selecting more than $\rho$ columns from
$C(X)$ gives zero, which proves the upper bound.  Put $A=C(v)$ and
choose an invertible real matrix $S=[U\ K]$ so that
\[
 S^{\mathsf T}AS=\begin{pmatrix} A_0&0\\0&0\end{pmatrix},
 \qquad \det A_0\ne0.
\]
For every $j$, the rank of $A+zC_j$ is at most $\rho$ for every real
$z$: every $(\rho+1)$-minor of the whole pencil is identically zero.
For sufficiently small $z$, its upper-left $\rho\times\rho$ block
remains invertible.  The Schur complement of this block must therefore
vanish identically.  Its coefficient of $z$ is exactly
$K^{\mathsf T}C_jK$, proving \eqref{deg:isotropic}.

The coefficient of $z^\rho$ in $\det(D+zA)$ is
\begin{equation}\label{deg:leading-direction}
 \frac{\det A_0}{(\det S)^2}\det(K^{\mathsf T}DK).
\end{equation}
This follows by taking the $zA_0$ entry in each of the first $\rho$
rows and columns after the congruence by $S$; the remaining block is
$K^{\mathsf T}DK$.  Thus a nonzero compression proves degree $\rho$.
Conversely, if the degree-$\rho$ homogeneous part is nonzero, it does
not vanish at a generic real maximal-rank direction, and
\eqref{deg:leading-direction} makes that compression nonsingular.
The conventions for $q=0$ are that $K$ has no columns and its
compressed determinant is one.
\end{proof}

There is a direct symbolic version that avoids expanding the full
$r\times r$ determinant.  Choose a nonsingular principal
$\rho\times\rho$ submatrix of $C(X)$ and, after a simultaneous
permutation of rows and columns, write
\[
 C(X)=\begin{pmatrix} A&B\\B^{\mathsf T}&F\end{pmatrix}.
\]
Over $\mathbb R(X)$, the rank condition gives $F=B^{\mathsf T}A^{-1}B$.
Set
\begin{equation}\label{deg:generic-kernel}
 K(X)=\begin{pmatrix}-A^{-1}B\\I_q\end{pmatrix}.
\end{equation}
Then the highest homogeneous part has the exact expression
\begin{equation}\label{deg:highest}
 [P]_{\rho}(X)=\det A\,
                  \det\bigl(K(X)^{\mathsf T}DK(X)\bigr).
\end{equation}
Although the displayed factors are rational functions, their product is
the indicated homogeneous polynomial.  The expression remains valid
when it vanishes.  A symmetric matrix of rank $\rho$ over a field of
characteristic zero has a nonsingular principal minor of size $\rho$;
this follows, for example, by symmetric elimination with $1\times1$
and $2\times2$ nonsingular pivots.

\subsection{Sharpness for moment matrices, including signed weights}

Let $w_n\in\mathbb R$ be a sequence and put
\begin{equation}\label{deg:moment}
 M_N=\sum_{n=0}^N w_n v(n)v(n)^{\mathsf T},
 \qquad v(n)=(1,n,\ldots,n^{r-1})^{\mathsf T}.
\end{equation}
Suppose that a decomposition
\begin{equation}\label{deg:decomp}
 M_N=D_N+\sum_{j=1}^t h_j(N)C_j
\end{equation}
holds with real $h_j(N)$ and real symmetric matrices $D_N,C_j$,
where the $C_j$ are independent of $N$.  Define
$P_N(X)=\det(D_N+C(X))$ and let $\delta_N$ be its total degree,
with degree $-\infty$ assigned to the zero polynomial.

\begin{theorem}[Only finitely many deficient cutoffs]
\label{deg:main}
Assume there are an integer $N_*\ge0$ and $\epsilon\in\{1,-1\}$ such
that $\epsilon w_n>0$ for every integer $n>N_*$.  Then
\[
 \delta_N\le\rho\quad\text{for all }N,
\]
and at most $r(r-\rho)$ integers $N\ge N_*$ satisfy
$\delta_N<\rho$.
If all nonzero $w_n$ have the same sign and at least $r$ distinct
nodes of positive absolute weight have been included by cutoff $N$,
then $\delta_N=\rho$ at that cutoff.

If in addition $D_N$ has entries rational in $N$, then
\begin{equation}\label{deg:generic-sharp}
 \deg_X\det(D_N+C(X))=\rho
 \qquad\text{over }\mathbb R(N).
\end{equation}
\end{theorem}
\begin{proof}
Choose one real maximal-rank direction and its kernel matrix $K$ as in
Lemma~\ref{deg:compression}.  Such a real direction exists because a
nonzero maximal minor is a nonzero real polynomial.  Equation
\eqref{deg:isotropic} and \eqref{deg:decomp} imply
\begin{equation}\label{deg:gram}
 G_N:=K^{\mathsf T}D_NK
 =K^{\mathsf T}M_NK
 =\sum_{n=0}^N w_n p(n)p(n)^{\mathsf T},
 \qquad p(n)=K^{\mathsf T}v(n).
\end{equation}
The components of $p$ are linearly independent real polynomials of
degree at most $r-1$.

For $N\ge N_*$, the matrix
\[
 \epsilon(G_{N+r}-G_N)
 =\sum_{n=N+1}^{N+r}\epsilon w_n p(n)p(n)^{\mathsf T}
\]
is positive definite.  Indeed, a nonzero linear combination of the
components of $p$ is a nonzero polynomial of degree at most $r-1$;
it cannot vanish at all $r$ displayed nodes.  Along any one of the
$r$ residue classes of cutoffs modulo $r$, the ordered eigenvalues of
$\epsilon G_N$ therefore increase strictly.  Each of the $q=r-\rho$
eigenvalues can equal zero at most once.  Consequently at most $q$
cutoffs in each residue class have singular $G_N$, and at most $rq$
cutoffs in total do.  By Lemma~\ref{deg:compression}, degree deficiency
can occur only at those cutoffs.  This also treats $q=0$, for which
there are no exceptions.

If the weights in the cutoff have one sign and at least $r$ distinct
nonzero nodes, the same polynomial argument makes $G_N$ definite
directly.  Finally, in the rational case $\det G_N$ is a rational
function of $N$.  It is nonzero at some integer cutoff by the preceding
argument, hence is not the zero rational function.  Formula
\eqref{deg:leading-direction} proves \eqref{deg:generic-sharp}.
\end{proof}

\begin{corollary}[Resolution for real rational weights]
\label{deg:rational}
Let $R$ be a nonzero real rational function having no pole at a
nonnegative integer.  Form its finite moment matrices and their
harmonic decomposition as in the finite-difference determinant report.
Then their formal harmonic degree over $\mathbb R(N)$ is exactly the
generic harmonic-pencil rank.  At most finitely many nonnegative
integer cutoffs have smaller degree.  For $R(n)>0$ at all
nonnegative integers, every $N\ge r-1$ has the exact generic degree.
\end{corollary}
\begin{proof}
A nonzero real rational function has constant nonzero sign for all
sufficiently large positive arguments.  The residue construction gives
the decomposition \eqref{deg:decomp} with rational $D_N$.  Apply
Theorem~\ref{deg:main}.  When poles and residues are represented using
nonreal conjugate pairs, replace each conjugate pair of formal harmonic
coordinates by its formal real and imaginary parts.  This invertible
linear change makes the coefficient matrices real and preserves both
generic rank and total degree.
\end{proof}

The real hypothesis in this proof is substantive: it supplies definiteness
of the sufficiently late increments.  No corresponding theorem for an
arbitrary complex rational weight is asserted here.  Likewise, formal
degree does not by itself prove a minimal multiple-polylogarithm depth
under functional relations.

\subsection{An exact test for exceptional cutoffs}

In the generic kernel formula \eqref{deg:generic-kernel}, replace $K$ by
the polynomial matrix
\[
 \widetilde K(X)=
 \begin{pmatrix}-\operatorname{adj}(A)B\\(\det A)I_q\end{pmatrix}.
\]
After clearing denominators in $N$, compute the polynomial
\begin{equation}\label{deg:exception-certificate}
 Q(N,X)=\det\bigl(\widetilde K(X)^{\mathsf T}D_N
                                      \widetilde K(X)\bigr).
\end{equation}
For an admissible cutoff $N_0$, degree is deficient if and only if
$Q(N_0,X)$ is identically zero.  Equivalently, all its coefficients in
$X$ vanish.  Over an exact coefficient field, their greatest common
divisor as polynomials in $N$ gives precisely the exceptional cutoffs
among its admissible nonnegative integer roots.  Denominator zeros must
be excluded or checked from the original finite moment formula.  This
uses a determinant of size $r-\rho$, not an expansion of the original
determinant.  A single fixed maximal-rank direction gives a cheaper
finite candidate list, which the generic test can then filter.

\subsection{A positive example and a genuine signed exception}

Take $r=3$ and first let
\[
 R(x)=\frac1{x+1}+\frac1{x+2}.
\]
There is one harmonic variable $X=H_N^{(1)}(1)$, and its coefficient is
\[
 C=\begin{pmatrix}2&-3&5\\-3&5&-9\\5&-9&17\end{pmatrix},
 \qquad \operatorname{rank}C=2,
 \qquad K=(2,3,1)^{\mathsf T}.
\]
Since $K^{\mathsf T}v(n)=(n+1)(n+2)$, the leading coefficient is
\begin{equation}\label{deg:positive-example}
 [P_N]_2=\frac{(N+1)(N+2)^2(N+3)}2 X^2.
\end{equation}
Here the top-left $2\times2$ minor of $C$ has determinant one.

For the signed weight
\begin{equation}\label{deg:signed-weight}
 R(x)=\frac1{x+1}-\frac7{5(x+2)}
     =\frac{3-2x}{5(x+1)(x+2)},
\end{equation}
the same kernel has compressed moment
\[
 K^{\mathsf T}D_NK
 =-\frac{(N-2)(N+1)(N+2)(N+3)}{10}.
\]
The relevant principal minor of $C$ has determinant $-7/5$, so
\begin{equation}\label{deg:signed-leading}
 [P_N]_2=\frac{7(N-2)(N+1)(N+2)(N+3)}{50}X^2.
\end{equation}
The only exceptional nonnegative cutoff is $N=2$, where the entire
formal polynomial is
\begin{equation}\label{deg:signed-exception}
 P_2(X)=\frac{620X-1137}{500}.
\end{equation}
Substituting the actual $X=H_2^{(1)}(1)=11/6$ gives $-1/1500$, also
equal to the three-node Vandermonde square times
$R(0)R(1)R(2)$.  Thus a fixed-cutoff exception is real, even when the
finite moment determinant itself is nonzero.  It does not contradict
generic sharpness over $\mathbb R(N)$.

% END sections/05_degree_rank.tex

% BEGIN sections/06_spectral_hierarchy.tex
\section{A critical finite-difference hierarchy of every harmonic rank}
\label{hier:sec}

The original finite-difference construction uses a stencil of order
$2r-2$ and confines its special-function dependence to one moment entry.
Shorter stencils give an explicit hierarchy whose harmonic ranks run
through every value from one to $r$, while the integer-order weights
retain the same critical decay.  Its spectral continuation has a
completely determined staircase of pole orders.  This is a spectral
finite-difference family; the more general overlapping resonances of
the earlier Gauss ensemble are a separate question.

Fix integers
\begin{equation}\label{hier:parameters}
 r\ge2,\qquad 1\le\rho\le r,\qquad
 d=2r-2,\qquad h=d-\rho+1=2r-\rho-1.
\end{equation}
For $\Re a>0$, use the right-half-plane logarithm to define
\begin{equation}\label{hier:weight}
 w_n^{[h]}(s,a)=\frac1{h!}\sum_{j=0}^h(-1)^j\binom hj
                                      (n+a+j)^{-s}.
\end{equation}
The corresponding ordered ensemble is
\begin{equation}\label{hier:ensemble}
 Z_{r,\rho}(s,a)=
 \sum_{0\le n_1<\cdots<n_r}
 \Delta(n_1,\ldots,n_r)^2\prod_{i=1}^r w_{n_i}^{[h]}(s,a).
\end{equation}
There is no additional $1/r!$ because the indices are strictly ordered.

\subsection{Finite moments and exact shifted-zeta closure}

Define finite polynomials and endpoint functions by
\begin{align}
 b_k(a;h)&=\frac1{h!}\sum_{j=0}^h(-1)^j\binom hj(-a-j)^k,
                                                        \label{hier:b}\\
 E_k(s,a;h)&=-\frac1{h!}\sum_{j=1}^h(-1)^j\binom hj
       \sum_{\ell=0}^{j-1}(\ell-j)^k(\ell+a)^{-s}.
                                                        \label{hier:E}
\end{align}
The convention $0^0=1$ applies to zeroth moments.  Polynomial
annihilation gives
\begin{equation}\label{hier:annihilation}
 b_0=\cdots=b_{h-1}=0,\qquad b_h=1.
\end{equation}
For $0\le\ell\le\rho-1$ let $C_\ell(a;h)$ be the symmetric Hankel
matrix with entries
\begin{equation}\label{hier:C}
 (C_\ell)_{ij}=\binom{i+j}{\ell}b_{i+j-\ell}(a;h),
 \qquad 0\le i,j<r,
\end{equation}
with a term set to zero if its lower index is negative or if
$\ell>i+j$.  Also let $E(s,a;h)=(E_{i+j}(s,a;h))_{i,j<r}$.

\begin{theorem}[A finite shifted-zeta determinant]
\label{hier:closure}
The ensemble \eqref{hier:ensemble} is absolutely and locally normally
convergent for $\Re s>\rho$, $\Re a>0$, and there
\begin{equation}\label{hier:closed}
 \boxed{\displaystyle
 Z_{r,\rho}(s,a)=\det\left(
 E(s,a;h)+\sum_{\ell=0}^{\rho-1}
                    \zeta(s-\ell,a)C_\ell(a;h)\right).}
\end{equation}
The right side gives its meromorphic continuation to all $s\in\mathbb C$.
Its only possible poles are $1,2,\ldots,\rho$.
\end{theorem}
\begin{proof}
For $\Re s>0$ the weight has the integral representation
\begin{equation}\label{hier:laplace}
 w_n^{[h]}(s,a)=\frac1{h!\Gamma(s)}
 \int_0^\infty t^{s-1}e^{-(n+a)t}(1-e^{-t})^h\,dt.
\end{equation}
Repeated finite differences, or expansion of this kernel at zero,
give locally uniformly
\begin{equation}\label{hier:decay}
 w_n^{[h]}(s,a)=\frac{(s)_h}{h!}(n+a)^{-s-h}
                     +O((n+1)^{-\Re s-h-1}).
\end{equation}
For ordered indices, $|\Delta(\boldsymbol n)|^2\le
\prod_i(n_i+1)^{2(i-1)}$.  The least decaying factor is summable
when $\Re s+h-d>1$, precisely $\Re s>\rho$.

For a moment of degree $k\le d$, initially take $\Re s>d+1$ so
that all separately shifted sums converge.  Shift the index by $j$
and expand
\[
 (n-j)^k=\sum_{\ell=0}^k\binom k\ell
                      (-a-j)^{k-\ell}(n+a)^\ell.
\]
The lower endpoints give \eqref{hier:E}, and the infinite shifted
sums give
\begin{equation}\label{hier:moment}
 \sum_{n\ge0}n^k w_n^{[h]}(s,a)
 =E_k(s,a;h)+\sum_{\ell=0}^{k-h}
     \binom k\ell b_{k-\ell}(a;h)\zeta(s-\ell,a).
\end{equation}
An empty sum is zero.  Both sides are holomorphic on their common
normal-convergence half-plane, so the identity theorem extends this
identity to $\Re s>\rho$.  Finite Cauchy--Binet and passage to the
normal limit give \eqref{hier:closed}.  Each endpoint function is
entire in $s$, and every Hurwitz factor has just its displayed pole.
\end{proof}

For later cutoff use, define
$H_N^{(z)}(a)=\sum_{n=0}^N(n+a)^{-z}$ for every complex $z$ and
\begin{equation}\label{hier:upper}
 T_{k,N}(s,a;h)=\frac1{h!}\sum_{j=1}^h(-1)^j\binom hj
       \sum_{\ell=N+1}^{N+j}(\ell-j)^k(\ell+a)^{-s}.
\end{equation}
The exact finite counterpart of \eqref{hier:moment} is
\begin{equation}\label{hier:finite-moment}
 \sum_{n=0}^N n^k w_n^{[h]}(s,a)
 = E_k+T_{k,N}+
   \sum_{\ell=0}^{k-h}\binom k\ell b_{k-\ell}
                                     H_N^{(s-\ell)}(a).
\end{equation}
It holds for every $s$, including orders at which the infinite
ensemble has a pole.  It follows from the finite index-shift identity;
no divergent series is split in its proof.

\subsection{Exact rank and the unique highest harmonic monomial}

\begin{theorem}[Every rank at the same critical decay]
\label{hier:rank}
At $s=\rho$ and real $a>0$, the rational weight
$R_{r,\rho,a}(n)=w_n^{[h]}(\rho,a)$ is positive and obeys
\begin{equation}\label{hier:critical}
 R_{r,\rho,a}(n)=\binom{2r-2}{\rho-1}n^{-(2r-1)}
                                      +O(n^{-2r}).
\end{equation}
Its harmonic pencil has generic rank exactly $\rho$.

Put $L=r-\rho$ and let
\begin{equation}\label{hier:Vcutoff}
 V_{L,N}(a)=\det\left(
   \sum_{n=0}^N n^{i+j}R_{r,\rho,a}(n)\right)_{0\le i,j<L},
 \qquad V_{0,N}=1.
\end{equation}
This is rational in $N$ (and rational in $a$).  The degree-$\rho$
homogeneous part of the full $r$-particle cutoff determinant, with
formal harmonic variables $X_q$ corresponding to $H_N^{(q)}(a)$, is
exactly
\begin{equation}\label{hier:top}
 \boxed{\displaystyle
 (-1)^{\rho(\rho-1)/2}V_{L,N}(a)X_\rho^\rho.}
\end{equation}
It is nonzero whenever $N\ge L-1$, with this cutoff condition vacuous
when $L=0$.
\end{theorem}
\begin{proof}
Positivity follows from \eqref{hier:laplace}.  In
\eqref{hier:decay}, $h+\rho=2r-1$ and
$(\rho)_h/h!=\binom{2r-2}{\rho-1}$.

By \eqref{hier:annihilation}, $C_\ell$ is supported in its final
$\rho-\ell$ rows and columns.  In its final
$(\rho-\ell)\times(\rho-\ell)$ block, every entry below the first
possible anti-diagonal in moment order vanishes, and the entries on
that anti-diagonal all equal
$\binom{h+\ell}{\ell}$.  In particular, all matrices lie in the
same final $\rho\times\rho$ block.  Within that block $C_0$ has
zeros for relative indices $i+j<\rho-1$ and ones when
$i+j=\rho-1$; all other $C_\ell$ start strictly after that
anti-diagonal.  Hence
\begin{equation}\label{hier:block-det}
 \det\left(\sum_{\ell=0}^{\rho-1}Y_\ell
                  (C_\ell)_{\rm bottom}\right)
     =(-1)^{\rho(\rho-1)/2}Y_0^\rho.
\end{equation}
Only the reversing permutation can contribute: its column-index sum
already equals the minimum possible sum of nonzero entries in all
rows.  Thus the pencil has rank $\rho$.

At $s=\rho$, $Y_\ell$ is the formal variable $X_{\rho-\ell}$.
To obtain degree $\rho$ in the full determinant one must take all
these entries from the final block.  The complementary determinant
is \eqref{hier:Vcutoff}, which proves \eqref{hier:top}.
Its moments have degree at most $2L-2<h$, so
\eqref{hier:finite-moment} expresses each using finite rational
endpoints alone.  Positivity of the weight makes this $L\times L$
Gram matrix positive definite once at least $L$ distinct nodes have
been included.  This proves the last assertion, including the empty
case $L=0$.
\end{proof}

For example, $r=3,\rho=2,h=3$ gives the positive critical weight
\[
 R(x)=\frac16\left\{\frac1{(x+a)^2}-\frac3{(x+a+1)^2}
              +\frac3{(x+a+2)^2}-\frac1{(x+a+3)^2}\right\}.
\]
Its two coefficient matrices, indexed by harmonic orders, are
\[
 C^{(1)}=\begin{pmatrix}0&0&0\\0&0&0\\0&0&4\end{pmatrix},
 \qquad
 C^{(2)}=\begin{pmatrix}0&0&0\\0&0&1\\0&1&-4a-6\end{pmatrix}.
\]
The exact leading harmonic term is
$-\bigl(\sum_{n=0}^N R(n)\bigr)X_2^2$, and its coefficient is
nonzero for every $N\ge0$.

\subsection{Exact pole orders and the leading Laurent coefficients}

For an integer $m$ with $1\le m\le\rho$, put
\begin{equation}\label{hier:pole-parameters}
 \eta=\rho-m+1,\qquad L=r-\eta,
 \qquad \beta_m=\binom{h+m-1}{m-1}.
\end{equation}
Define the convergent lower-particle normalization
\begin{equation}\label{hier:lower-normalization}
 G_{L,h}^{(m)}(a)=
 \det\left(\sum_{n\ge0}n^{i+j}w_n^{[h]}(m,a)\right)_{0\le i,j<L},
 \qquad G_{0,h}^{(m)}=1.
\end{equation}

\begin{theorem}[A staircase of exact spectral poles]
\label{hier:poles}
For real $a>0$, $Z_{r,\rho}(s,a)$ has a pole of exact order
$\rho-m+1$ at every $m=1,\ldots,\rho$.  Its leading coefficient is
\begin{equation}\label{hier:leading-Laurent}
 \boxed{\displaystyle
 \lim_{s\to m}(s-m)^{\eta}Z_{r,\rho}(s,a)
 =(-1)^{\eta(\eta-1)/2}\beta_m^{\eta}G_{L,h}^{(m)}(a).}
\end{equation}
The lower normalization is strictly positive.  For complex $a$ with
$\Re a>0$, the same identity holds holomorphically, but at special
complex parameters the leading coefficient is not asserted nonzero.
\end{theorem}
\begin{proof}
Only $\zeta(s-m+1,a)$ has a pole at $s=m$; its residue is one.
Consequently the continued moment matrix has the local form
\[
 M(s)=\frac{C_{m-1}}{s-m}+A(s)
\]
with $A$ holomorphic.  The residue matrix is supported in its final
$\eta\times\eta$ block.  This block has zeros before the
anti-diagonal and constant anti-diagonal $\beta_m$, so its
determinant is
$(-1)^{\eta(\eta-1)/2}\beta_m^{\eta}$.

In the full determinant, the coefficient of $(s-m)^{-\eta}$ is this
residue determinant times the complementary top-left $L\times L$
minor of $A(m)$.  These top-left entries have no pole and coincide
with ordinary moments of $w_n^{[h]}(m,a)$.  Indeed, their largest
degree is $2L-2$ and the weight decays as $n^{-m-h}$; the required
strict inequality is
\[
 m+h-(2L-2)>1,
\]
which here has left side $\eta+2>1$.  The same inequality holds in
a small complex neighborhood, so normal convergence identifies the
analytic moment entries with their series there.  This proves
\eqref{hier:leading-Laurent}.  At positive real $a$, the Laplace
formula makes the weight strictly positive on infinitely many nodes,
so the lower Gram determinant is strictly positive.  The pole order
is therefore exact.
\end{proof}

\subsection{A Gamma product at the highest-order pole}

At $m=1$, put $L=r-\rho$.  The normalization in the preceding theorem
has the finite product
\begin{equation}\label{hier:Gamma-product}
 \boxed{\displaystyle
 G_{L,h}^{(1)}(a)=
 \prod_{j=0}^{L-1}
 \frac{(j!)^2\Gamma(a+j)\Gamma(h-L+1-j)}
      {\Gamma(h+1-j)\Gamma(a+h-j)}.}
\end{equation}
In particular, the leading Laurent coefficient at $s=1$ is
$(-1)^{\rho(\rho-1)/2}$ times this product.  For $\rho=r$ it is
the nonzero constant $(-1)^{r(r-1)/2}$.

Here is a derivation that also fixes the normalization.  The beta
partial-fraction identity gives
\[
 w_n^{[h]}(1,a)=\frac{\Gamma(n+a)}{\Gamma(n+a+h+1)}.
\]
Replace the two monomial bases of the $L$-particle moment determinant
by the monic bases $n^{\underline i}$ and $(n+a)_j$.  The mixed
moment is
\begin{equation}\label{hier:mixed-moment}
 M_{ij}=\frac{\Gamma(i+1)\Gamma(a+i+j)\Gamma(h-i-j)}
                  {\Gamma(h+1-j)\Gamma(a+h)},
 \qquad 0\le i,j<L.
\end{equation}
To see this, write $n=i+k$ and sum the resulting Gauss series with
parameters $i+1,a+i+j,a+i+h+1$.  Its excess is $h-i-j>0$;
in our parameters $h>2L-2$.  Thus the sum is justified by absolute
convergence.

After factoring rows and columns in \eqref{hier:mixed-moment}, the
remaining determinant is
\[
 \det_{0\le i,j<L}\frac{(a+i)_j}{(1-h+i)_j}.
\]
The elementary factorial determinant
\[
 \det\frac{(x_i)_j}{(x_i+t)_j}
 =\frac{\prod_{j=0}^{L-1}(t)_j\,
             \prod_{i<k}(x_k-x_i)}
            {\prod_{i=0}^{L-1}(x_i+t)_{L-1}}
\]
follows by multiplying each row by $(x_i+t)_{L-1}$ and comparing
polynomial coefficients; evaluation at $x_i=-i$ computes the
coefficient determinant.  Apply it with $x_i=a+i$ and $t=1-h-a$.
Use
\[
 (1-h-a)_j=(-1)^j\frac{\Gamma(a+h)}{\Gamma(a+h-j)},\qquad
 (1-h+i)_{L-1}=(-1)^{L-1}
                    \frac{\Gamma(h-i)}{\Gamma(h-i-L+1)}.
\]
All signs cancel, and the remaining factors give
\eqref{hier:Gamma-product}.  This is the same classical factorial
determinant mechanism as the predecessor's residue product; the new
point is its role as the coefficient of a pole of arbitrary order
$\rho$ in the shorter-stencil spectral family.

\subsection{A complete two-particle formula and its Stieltjes finite parts}

For $r=\rho=2$ one has $h=1$.  Directly from the three moment
identities,
\begin{equation}\label{hier:r2}
 \boxed{\displaystyle
 Z_{2,2}(s,a)=2a^{-s}\zeta(s-1,a)
        +(1-2a)a^{-s}\zeta(s,a)-\zeta(s,a)^2.}
\end{equation}
Thus, for every integer $m\ge3$, the ordinary convergent sum
\begin{align}\label{hier:r2-sum}
 &\sum_{0\le n<k}(k-n)^2
 \bigl((n+a)^{-m}-(n+a+1)^{-m}\bigr)
 \bigl((k+a)^{-m}-(k+a+1)^{-m}\bigr)\notag\\
 &\hspace{12mm}=2a^{-m}\zeta(m-1,a)
            +(1-2a)a^{-m}\zeta(m,a)-\zeta(m,a)^2.
\end{align}
For example, $a=1,m=3$ gives
$2\zeta(2)-\zeta(3)-\zeta(3)^2$.

Use the convention
\[
 \zeta(1+u,a)=\frac1u+
    \sum_{k\ge0}\frac{(-1)^k\gamma_k(a)}{k!}u^k,
 \qquad c=\frac1a-2.
\]
Expanding \eqref{hier:r2} at the double pole gives the full first
three Laurent terms
\begin{align}\label{hier:r2-s1}
 Z_{2,2}(1+u,a)
 ={}&-\frac1{u^2}+\frac{c-2\gamma_0(a)}u\notag\\
 &+2\gamma_1(a)-\gamma_0(a)^2
       +c\bigl(\gamma_0(a)-\log a+1\bigr)+O(u).
\end{align}
At the simple pole $s=2$,
\begin{align}\label{hier:r2-s2}
 \operatorname{FP}_{s=2}Z_{2,2}(s,a)
 ={}&2a^{-2}\bigl(\gamma_0(a)-\log a\bigr)
 +(1-2a)a^{-2}\zeta(2,a)-\zeta(2,a)^2.
\end{align}
These may be written entirely in generalized Stieltjes and polygamma
coordinates using $\gamma_0(a)=-\psi(a)$ and $\zeta(2,a)=\psi'(a)$.
They are spectral Laurent finite parts.  No equality with an
unsubtracted endpoint sum or a hard-cutoff constant is implicit.

There is a particularly simple check at order zero. Substitution of
$\zeta(0,a)=1/2-a$ and
$\zeta(-1,a)=-(a^2-a+1/6)/2$ into \eqref{hier:r2} gives
\begin{equation}\label{hier:zero-contact}
                  Z_{2,2}(0,a)=\frac1{12}
\end{equation}
for its spectral continuation. Yet at $s=0$ every weight
$(n+a)^0-(n+a+1)^0$ is exactly zero, so the unmodified defining
sum at that parameter has ordinary value zero. The discrepancy is
an explicit example of noncommutation between continuation and
specialization of the original infinite sum; \eqref{hier:zero-contact}
does not assign $1/12$ to an ordinary sum of zero terms.

All further Laurent coefficients of \eqref{hier:closed} are obtained
by finite determinant arithmetic from the Hurwitz expansions at the
relevant shifted orders.  In particular, at $s=m$ the coordinates are
the generalized Stieltjes constants from the sole resonant Hurwitz
factor, together with finitely many ordinary Hurwitz derivatives from
the nonresonant factors and elementary endpoint logarithms.  This is
an exact coefficient-generation procedure, without a claim that the
resulting higher constants admit an additional classical reduction.

% END sections/06_spectral_hierarchy.tex

% BEGIN sections/07_polylog_generator.tex
\section{A cutoff polylogarithm identity and its Euler primitive}
\label{hier:gf-section}

The first rank-two member of the new hierarchy gives a compact
polylogarithmic consequence.  Its weight is
\begin{equation}\label{hier:gf-weight}
 R(n)=\frac1{(n+1)^2}-\frac1{(n+2)^2},
\end{equation}
whereas the predecessor's explicit repeated-pole generator uses the
distinct second-difference weight
$\tfrac12\{(n+1)^{-2}-2(n+2)^{-2}+(n+3)^{-2}\}$.
Define the finite ordered determinant sum
\begin{equation}\label{hier:gf-cutoff}
 Z_N=\sum_{0\le n<k\le N}(k-n)^2R(n)R(k),\qquad Z_0=0.
\end{equation}

Our one-direction multiple-polylogarithm convention is
\begin{equation}\label{hier:gf-convention}
 \operatorname{Li}_{s_1,\ldots,s_j}(z)
 =\sum_{n_1>\cdots>n_j\ge1}
       \frac{z^{n_1}}{n_1^{s_1}\cdots n_j^{s_j}},
 \qquad |z|<1,
\end{equation}
with positive integer indices.  Thus only the largest summation index
carries $z$; in a multiple-argument convention this is
$\operatorname{Li}_{s_1,\ldots,s_j}(z,1,\ldots,1)$.
In particular the shuffle identity has normalization
$\operatorname{Li}_1(z)^2=2\operatorname{Li}_{1,1}(z)$.

\begin{theorem}[Finite cutoff and its complete generator]
\label{hier:gf-theorem}
For $M=N+1\ge1$,
\begin{equation}\label{hier:gf-finite}
 \boxed{\displaystyle
 Z_{M-1}=
 2\left(1-\frac1{(M+1)^2}\right)H_M
 -(H_M^{(2)})^2
 -\frac{M^2}{(M+1)^2}\bigl(H_M^{(2)}+1\bigr).}
\end{equation}
The generating function $G(z)=\sum_{N\ge0}Z_Nz^{N+1}$ satisfies
\begin{align}\label{hier:gf-closed}
 z(1-z)G(z)={}&2\operatorname{Li}_1(z)-\operatorname{Li}_2(z)
                 -z-z\operatorname{Li}_4(z)\notag\\
 &+2(1-z)\bigl(\operatorname{Li}_{1,2}(z)
                         -\operatorname{Li}_{2,1}(z)\bigr)
 -(1+z)\operatorname{Li}_{2,2}(z).
\end{align}
All series are absolutely and locally normally convergent for $|z|<1$.
The quotient in \eqref{hier:gf-closed} is interpreted by its removable
value at zero; $G(z)=\tfrac5{48}z^2+O(z^3)$.
\end{theorem}
\begin{proof}
Use the monic moment basis $1,n+1$, which leaves the determinant
unchanged, and put $p=M+1$.  Telescoping gives its three moments as
\[
 \mu_0=1-p^{-2},\qquad
 \mu_1=H_M^{(2)}-Mp^{-2},\qquad
 \mu_2=2H_M-H_M^{(2)}-1+2p^{-1}-p^{-2}.
\]
Expanding $\mu_0\mu_2-\mu_1^2$ proves \eqref{hier:gf-finite}.
The basic absolutely convergent summation identities are
\begin{align*}
 \sum_{M\ge1}H_M^{(p)}z^M
    &=\frac{\operatorname{Li}_p(z)}{1-z},\qquad
 \sum_{M\ge1}\frac{z^M H_M^{(p)}}{(M+1)^q}
    =z^{-1}\operatorname{Li}_{q,p}(z),\\
 \sum_{M\ge1}(H_M^{(2)})^2z^M
    &=\frac{2\operatorname{Li}_{2,2}(z)+\operatorname{Li}_4(z)}{1-z},
 \qquad
 \sum_{M\ge1}\frac{z^M}{(M+1)^q}
    =z^{-1}\operatorname{Li}_q(z)-1.
\end{align*}
The square identity follows by splitting its two inner indices into
the two strict orderings and the diagonal.  Substitute these formulas
into \eqref{hier:gf-finite}, after writing
$M^2/(M+1)^2=1-2/(M+1)+1/(M+1)^2$.  Collecting terms gives
\eqref{hier:gf-closed}.

Equation \eqref{hier:gf-finite} gives $Z_N=O(\log(N+2))$.
Every multiple-polylogarithm series in \eqref{hier:gf-convention}
also converges normally on compact subdisks, since its inner sums
grow at most polynomially while $z^{n_1}$ decays geometrically.
These estimates justify all interchanges.  The original generating
series proves that the apparent singularity at zero is removable;
$Z_0=0$ and $Z_1=R(0)R(1)=5/48$ give the stated initial terms.
\end{proof}

\begin{corollary}[An anchored Euler primitive]
\label{hier:gf-primitive}
For $|z|<1$, set
$I(z)=\int_0^zG(t)\,dt/t$ along any path in the disk.  Then
\begin{align}\label{hier:gf-primitive-closed}
 I(z)={}&\frac{
 \operatorname{Li}_1(z)^2+2\operatorname{Li}_{2,1}(z)
 -\operatorname{Li}_{1,2}(z)+\operatorname{Li}_{2,2}(z)
 +\operatorname{Li}_2(z)-\operatorname{Li}_1(z)}{z}\notag\\
 &-\operatorname{Li}_5(z)-\operatorname{Li}_{1,4}(z)
     -2\operatorname{Li}_{3,2}(z)
     -2\operatorname{Li}_{1,2,2}(z).
\end{align}
The removable value at zero is $I(0)=0$, and
\[
 I(z)=\sum_{N\ge0}\frac{Z_N}{N+1}z^{N+1}
     =\frac5{96}z^2+O(z^3).
\]
\end{corollary}
\begin{proof}
Let $\theta=z\,d/dz$.  From the absolutely convergent series,
\[
 \theta\operatorname{Li}_{p,\boldsymbol s}
      =\operatorname{Li}_{p-1,\boldsymbol s}\quad(p>1),\qquad
 \theta\operatorname{Li}_{1,\boldsymbol s}
      =\frac{z}{1-z}\operatorname{Li}_{\boldsymbol s}.
\]
For an empty $\boldsymbol s$ the second formula means
$\theta\operatorname{Li}_1=z/(1-z)$.
These identities also prove
$\operatorname{Li}_1^2=2\operatorname{Li}_{1,1}$ by comparing
derivatives and values at zero.  Apply them to
\eqref{hier:gf-primitive-closed}: its Euler derivative is the function $G(z)$ specified by
\eqref{hier:gf-closed}.  The convergent expansions at zero give the
zero integration constant and the stated removability.  Alternatively,
coefficientwise division of \eqref{hier:gf-finite} by $M$ gives the
same formula by
\[
 \frac1{M(M+1)^2}
  =\frac1M-\frac1{M+1}-\frac1{(M+1)^2}.
\]
\end{proof}

The raw generator has a representation of depth at most two, although
its coefficient determinant has harmonic degree two.  Its displayed
Euler primitive uses depth at most three.  These are explicit
representations, not claims of minimal depth under functional identities.

\subsection{A convergent second Euler endpoint}

\begin{corollary}[An ordinary-zeta endpoint]
\label{hier:gf-endpoint}
The second Euler primitive has the absolutely convergent value
\begin{align}\label{hier:gf-endpoint-value}
 \sum_{N\ge0}\frac{Z_N}{(N+1)^2}
 &=\int_0^1\frac{-\log z}{z}G(z)\,dz\notag\\
 &=1+3\zeta(2)-2\zeta(3)-\frac34\zeta(4)
                     -\zeta(3)^2-\frac{19}{24}\zeta(6).
\end{align}
The new content here is the determinant endpoint identity; the
weight-six multiple-zeta reductions used below are classical.
\end{corollary}
\begin{proof}
The estimate $Z_N=O(\log(N+2))$ proves convergence.  Positivity of
$Z_N$ also permits Tonelli's theorem in the integral representation.
Divide \eqref{hier:gf-finite} by $M^2$ and collect its rational
coefficients to obtain the pointwise identity
\[
 \frac{Z_{M-1}}{M^2}
 =-\frac{(H_M^{(2)})^2}{M^2}
  -\frac{2H_M}{(M+1)^2}
  +\frac{4H_M}{M(M+1)}
  -\frac{H_M^{(2)}+1}{(M+1)^2}.
\]
Every term is separately summable.  The standard nested-series
interpretation and elementary telescoping give
\[
 \sum_{M\ge1}\frac{H_M}{M(M+1)}=\zeta(2),\qquad
 \sum_{M\ge1}\frac{H_M}{(M+1)^2}=\zeta(2,1)=\zeta(3),
 \qquad
 \sum_{M\ge1}\frac{H_M^{(2)}}{(M+1)^2}=\zeta(2,2)=\frac34\zeta(4).
\]
Splitting coincident and strictly ordered indices gives
\begin{equation}\label{hier:gf-square-Euler}
 \sum_{M\ge1}\frac{(H_M^{(2)})^2}{M^2}
 =2\zeta(2,2,2)+2\zeta(4,2)+\zeta(2,4)+\zeta(6)
 =\zeta(3)^2+\frac{19}{24}\zeta(6).
\end{equation}
For clarity, the three weight-six inputs are
\[
 \zeta(2,2,2)=\frac3{16}\zeta(6),\qquad
 \zeta(4,2)=\zeta(3)^2-\frac43\zeta(6),\qquad
 \zeta(2,4)=\frac{25}{12}\zeta(6)-\zeta(3)^2.
\]
They follow from short classical calculations.  Expanding the cube
of $\zeta(2)$ by coincident indices gives the first.  Euler's formula
$2\zeta(m,1)=m\zeta(m+1)-\sum_{j=1}^{m-2}
\zeta(j+1)\zeta(m-j)$, with $m=5$, gives
$\zeta(5,1)=\tfrac34\zeta(6)-\tfrac12\zeta(3)^2$;
see \cite[Eq.~(1.3)]{hier:BB2005}.  Equating the convergent shuffle
and stuffle products
\[
 \zeta(3)^2=2\zeta(3,3)+6\zeta(4,2)+12\zeta(5,1),
 \qquad
 \zeta(3)^2=2\zeta(3,3)+\zeta(6)
\]
then gives the second weight-six input.  Finally,
$\zeta(2,4)+\zeta(4,2)=\zeta(2)\zeta(4)-\zeta(6)$ gives the third.
The same Euler formula at $m=2$ supplies $\zeta(2,1)=\zeta(3)$.
Substitution into the summable pointwise identity proves
\eqref{hier:gf-endpoint-value}.
\end{proof}

% END sections/07_polylog_generator.tex

% BEGIN sections/08_order_ten.tex
% Research fragment: all labels begin ten:. Standard amsmath/amsthm required.
% Cite AlternatingIncoming, QuarticIncoming, and GeneralHankel as the inspected project sources.
\section{Complete reconstruction for smaller tail order ten}
\label{ten:section}

The two latest incoming reconstruction reports leave a particularly
small but genuine gap.  The alternating-tail report proves the order-ten
prefix theorem outside $b\equiv962\pmod{2048}$, while the quartic
reconstruction report proves it for every $b\ge1920$
\cite{AlternatingIncoming,QuarticIncoming}.  We complete the family.  The new step is
a coefficientwise congruence with one additional bit of precision:
its only exceptional class is $b\equiv3010\pmod{4096}$, which lies
entirely inside the range of the second report's real-variable argument.
We reproduce that argument at order ten, with explicit constants, to
make the resulting nonvanishing proof complete.  The small determinant
normalization and its Pad\'e origin are credited to the source reports;
the modulus-$2^{17}$ congruence and the completed family are the present
contribution.  No claim of arithmetic independence of numerical special
values is involved.

\subsection{Statement and exact meaning of reconstruction}

Let
\[
 A_s(n)=\zeta(s,n+1)=\zeta(s)-H_n^{(s)},\qquad x_n=n+\tfrac12,
 \qquad
 D_{10,2b}(w)=\sum_{n\ge0}A_{10}(n)A_{2b}(n)x_n^{-w}.
\]
The series initially converges for $\operatorname{Re}w>-2b-7$.
Its centered Euler--Maclaurin continuation is regular at nonpositive
even integers.  Write $D_r=D_{10,2b}(-2r)$ and use the source's
triangular coordinates
\begin{align}
 C_j&=\sum_{r=0}^{j}\binom{2j+1}{2r}2^{-2(j-r)}D_r,
 \label{ten:C}\\
 D_r&=\sum_{j=0}^{r}\frac{\binom{2r}{2j}}{2j+1}
             B_{2r-2j}(1/2)C_j.
 \label{ten:D}
\end{align}
For $j\ge b$, each $C_j$ is a rational combination of
$Z_k=\zeta(2k)$, $0\le k<b$, with $Z_0=-1/2$.  The words
\emph{formal rank} refer to rational coefficient vectors in the
formal basis
\[
 C_0,\ldots,C_{b-1},Z_0,\ldots,Z_{b-1};
\]
they do not refer to the linear independence of the values after
evaluation.

Put $\mu_j=B_{2j}/(2j)!$ for $j\ge0$, and $\mu_j=0$ for $j<0$.
The late coordinate matrix is
\begin{equation}\label{ten:M}
 (M_{5,b})_{tk}=\mu_{t+k-4}
 \left\{\frac{(2b+2t-9)!}{(2b-2k)!}
 +\mathbf1_{k<5}\frac{(2t+1)!}{(10-2k)!}\right\},
 \qquad 0\le t,k<b.
\end{equation}
In particular,
$(C_b,\ldots,C_{2b-1})^{\mathsf T}
=M_{5,b}(Z_0,\ldots,Z_{b-1})^{\mathsf T}$.
Formula \eqref{ten:M} also defines the late coordinate row for
every integer $t\ge0$.  Its binomial form, useful without factorial
conventions, is
\[
 B_{2t+2k-8}\left\{
 \frac{\binom{2t+2}{10-2k}}{2t+2}
 +\frac{\binom{2b+2t-8}{2b-2k}}{2b+2t-8}\right\},
\]
where negative Bernoulli indices give zero and a summand is omitted
unless its lower binomial index lies between $2$ and its upper index.
These are the inherited finite Bernoulli moment formulas
\cite{GeneralHankel}.  The full first-$2b$ coordinate determinant is
$\det M_{5,b}/(4b-1)!!$.

\begin{theorem}[Complete order-ten optimal prefix]\label{ten:main}
For every integer $b\ge6$, $M_{5,b}$ is nonsingular.  Consequently,
for every even $q>10$, the $q$ moments
\[
 D_{10,q}(0),D_{10,q}(-2),\ldots,D_{10,q}(-2(q-1))
\]
have full formal rank $q$.  Every subsequent centered even moment
is an explicit rational linear combination of this prefix.
\end{theorem}

The conclusion now holds for every unequal even pair whose smaller
order is at most ten, using the previously established smaller-order
families.  It does not establish the unrestricted conjecture at
smaller orders twelve and above, nor the full gap-six family.

\subsection{A finite five-dimensional definition of the numerator}

The following normalization is retained from \cite{AlternatingIncoming}.
For $0\le i,\ell\le4$, define
\[
 n_i=b-5+i,\quad N_{i\ell}=2(n_i-\ell)+1,
 \quad R_\ell=2b-8+2\ell,
\]
\[
 F_{i\ell}=\frac1{(2\ell+1)!}
       \prod_{r=0}^{2\ell-1}(2n_i-r)(2n_i+3+r),
\]
and define Newton coefficients $d_{ih}$ by the finite identity
\begin{equation}\label{ten:Newton}
 (2n_i+2-y)(4n_i+2-y)^{\underline{2i}}
       =\sum_{h=0}^{2i+1}d_{ih}y^{\underline h}.
\end{equation}
Equivalently,
\[
 d_{ih}=\frac1{h!}\sum_{r=0}^{h}(-1)^{h-r}\binom hr
       (2n_i+2-r)(4n_i+2-r)^{\underline{2i}}.
\]
For independent variables $U_1,\ldots,U_9$ put
\begin{equation}\label{ten:calA}
 \mathcal A_{i\ell}=2F_{i\ell}+
 \frac{R_\ell!}{(2n_i+1)(2n_i+2)(2\ell+2)!N_{i\ell}!}
 \sum_{k=0}^{i}d_{i,2k+1}N_{i\ell}^{\underline{2k+1}}
                 U_{5-i+\ell+k}.
\end{equation}
All factorial ratios with $b$-dependent arguments have fixed index
differences and therefore mean finite rational products in $b$.

Set
\begin{align*}
 \Delta_{10}(b)={}&4096b^{12}-110592b^{11}+1236992b^{10}
 -7303680b^9+23705088b^8\\
 &-38140416b^7+14925056b^6+24490080b^5-4469984b^4\\
 &-15838992b^3-2571948b^2+4244400b+1488375,
\end{align*}
and define the finite rational expression
\begin{equation}\label{ten:Pdef}
 P_{10}(b;U)=138240(2b-7)(2b-5)\Delta_{10}(b)
                       \frac{\det\mathcal A}{\det F}.
\end{equation}
Exact coefficient arithmetic gives
\begin{align}\label{ten:Fdet}
 \det F={}&\frac{16384}{4465125}(b-3)(b-2)^2(b-1)
     (2b-7)(2b-5)^2(2b-3)\notag\\
 &\cdot(4b-15)(4b-13)(4b-11)^2(4b-9)^2(4b-7)^2
          (4b-5)(4b-3)\Delta_{10}(b).
\end{align}
The expression $P_{10}$ belongs to $\mathbb Z[b,U_1,\ldots,U_9]$.
It has $439$ monomials when grouped by the $U$ exponents, or
$10403$ monomials after expanding also in $b$.  Its coefficients
have degree at most $24$ in $b$.  Equations
\eqref{ten:Newton}--\eqref{ten:Pdef} specify it completely, without
interpolation or a dependence on an unspecified fitted expression.
The accompanying exact verifier independently reconstructs
\eqref{ten:Fdet}, clears the five row denominators of
\eqref{ten:calA}, sums the $120$ signed permutation products, and
compares every coefficient of \eqref{ten:Pdef} with the inherited
integer certificate.  Thus the reused certificate is independently
replayed as an algebraic identity.

For $b\ge10$, with $H=(\mu_{t+k-4})_{0\le t,k<b}$, define
\[
 \rho_{5,b}=\frac{\det M_{5,b}}
 {\displaystyle
  \frac{\prod_{t=0}^{b-1}(2b+2t-9)!}
       {\prod_{k=0}^{b-1}(2b-2k)!}\det H}.
\]
The source reduction is
\begin{equation}\label{ten:rho}
 \rho_{5,b}
 =\left.\frac{\det\mathcal A}{\det F}\right|_{U_j=2B_{2b-2j}}
 =\frac{P_{10}(b;2B_{2b-2},\ldots,2B_{2b-18})}
 {138240(2b-7)(2b-5)\Delta_{10}(b)}.
\end{equation}
Here $H$ is nonsingular.  One direct proof conjugates $H$ by
$\operatorname{diag}(1,-1,1,\ldots)$ and partitions after its first
three rows and columns.  The resulting matrix has blocks
\[
 \begin{pmatrix}ee^{\mathsf T}&B\\B^{\mathsf T}&-P\end{pmatrix},
 \qquad e=(0,0,1)^{\mathsf T},
\]
where $P$ is positive definite by the representation
$2\sum_{m\ge1}(2\pi m)^{-2s}p((2\pi m)^{-2})^2>0$ for
$s\ge1$ and a nonzero polynomial $p$.  The first two rows of $B$
have a nonsingular antitriangular minor.  Consequently
$ee^{\mathsf T}+BP^{-1}B^{\mathsf T}$ is positive definite.

For completeness, the exact nonzero multiplier behind
\eqref{ten:rho} can be obtained without assuming Pad\'e normality.
The monic polynomials
\[
 \pi_n(x)=\sum_{t=0}^{n}
 \frac{2(2n+2t+1)!(2n+1)!}
 {(4n+2)!(2n-2t+1)!(2t)!}x^t
\]
satisfy, for $\mathcal N(x^r)=\mu_{r+1}$,
\[
 \mathcal N(x^j\pi_n)=0\ (j<n),\qquad
 \mathcal N(x^n\pi_n)=h_n,
 \quad h_n=\frac{((2n+1)!)^2}{(4n+2)!(4n+3)!}.
\]
These finite identities follow by coefficient comparison in the
diagonal Pad\'e formula for $e^t$, as in \cite{QuarticIncoming}; they
can also be checked directly using the Bernoulli generating function.
The same coefficients give
\[
 \sum_{t=0}^{n}[x^t]\pi_n(x)\,\mu_{t-\ell}
 =c_n f_\ell(n),\qquad
 c_n=\frac{(2n+1)(2n+2)!}{(4n+2)!},
\]
where $f_\ell(n)=F_{i\ell}$ when $n=n_i$.
Apply these monic row combinations to
$H+U^{\mathrm{pert}}$, where
\[
 U^{\mathrm{pert}}_{tk}
 =\mathbf1_{k<5}\mu_{t+k-4}
 \frac{(2t+1)!(2b-2k)!}{(10-2k)!(2b+2t-9)!}.
\]
The last five transformed rows vanish in columns $k\ge5$.
The complementary block is upper triangular with diagonal
$h_0,\ldots,h_{b-6}$.  Reverse the retained five columns and
divide row $i$ by $c_{n_i}$.  The unperturbed entries become $F$.
For the perturbation use the elementary Bernoulli identity
\[
 \sum_{j=0}^{N}\binom Nj B_j(N-j)^{\underline h}
       =N^{\underline h}B_{N-h}(1).
\]
Since $b\ge10$, all the retained indices $N-h$ are at least two;
only odd $h$ survive.  Restoring the omitted $j=1$ term supplies
one further copy of $F$, explaining the necessary factor $2F$
in \eqref{ten:calA}.  This proves
\begin{equation}\label{ten:multiplier}
 \det M_{5,b}=(-1)^{5b-15}
 \frac{\prod_{t=0}^{b-1}(2b+2t-9)!}
      {\prod_{k=0}^{b-1}(2b-2k)!}
 \left(\prod_{n=0}^{b-6}h_n\right)
 \left(\prod_{i=0}^{4}c_{n_i}\right)\det\mathcal A.
\end{equation}
The same argument with the perturbation absent gives the identical
multiplier for $\det H$ and $\det F$, proving \eqref{ten:rho}.

\subsection{One additional bit removes the remaining finite obstruction}

All congruences in this subsection are in $\mathbb Z_{(2)}$,
the rationals with odd denominator.  We use the classical fact
$v_2(B_{2m})=-1$ for $m\ge1$.

\begin{lemma}[A doubled Bernoulli congruence]\label{ten:B17}
For every integer $m\ge13$,
\[
 2B_{2m}\equiv\mathcal B_{17}(m)\pmod{131072},
\]
where
\begin{equation}\label{ten:Bpoly}
 \mathcal B_{17}(x)=1+20068x-26384x^2+2816x^3+45824x^4
             +4096x^5-24576x^6+65536x^8.
\end{equation}
\end{lemma}
\begin{proof}
Faulhaber's identity, with $L=512$, gives
\[
 \sum_{a=0}^{L-1}a^{2m}
 =LB_{2m}+\sum_{h=2}^{2m+1}
      \binom{2m}{h-1}\frac{L^h}{h}B_{2m+1-h}.
\]
The $h=2$ term vanishes.  Every term with $h\ge3$ has valuation
at least $9h-v_2(h)-1\ge26$.  Divide by $256$ and reduce modulo
$2^{17}$.  The even $a$ terms vanish because $2m\ge26$.
For odd $a$, expand
$a^{2m}=(1+(a^2-1))^m$; since $8\mid a^2-1$, terms with
binomial degree at least nine have valuation at least $27-8>17$
after division.  Therefore
\begin{equation}\label{ten:finite-average}
 2B_{2m}\equiv\frac1{256}
 \sum_{a=1,3,\ldots,511}\sum_{j=0}^{8}
          \binom mj(a^2-1)^j\pmod{2^{17}}.
\end{equation}
Expand the finite right side as a polynomial in $m$.  Its rational
coefficients all have odd denominators, and their reductions modulo
$2^{17}$, in increasing powers, are
\[
 (1,20068,-26384,2816,45824,4096,-24576,0,65536).
\]
This is an exact finite coefficient calculation, reproduced by the
verifier from the $256$ odd support points and the nine displayed
binomial terms.  It gives \eqref{ten:Bpoly} and proves the lemma.
\end{proof}

Define
\begin{equation}\label{ten:Rpoly}
 \begin{split}
 \mathcal R(b)={}&1536b^8+128b^7-928b^6-816b^5+464b^4\\
                 &-1696b^3-1930b^2-1831b+694.
 \end{split}
\end{equation}
The new coefficientwise polynomial identity is
\begin{equation}\label{ten:congruence}
 P_{10}\bigl(b;\mathcal B_{17}(b-1),\ldots,
                    \mathcal B_{17}(b-9)\bigr)
       \equiv32\mathcal R(b)\pmod{131072}.
\end{equation}
This is verified in $\mathbb Z[b]$, not only at selected integer
arguments.  A compact exact procedure is to group the numerator's
$10403$ terms by their $439$ $U$-monomials, substitute the nine
explicit polynomials, multiply modulo $131072$, and compare every
coefficient with \eqref{ten:Rpoly}.  The provided script carries
out precisely those finite ring operations, after independently
rebuilding the numerator from \eqref{ten:calA}.  Thus
\eqref{ten:congruence} is a finite algebraic certificate whose
entire input is displayed or defined by finite products above.

The linear coefficient of $\mathcal R$ is odd, and every coefficient
of degree at least two is even.  Hence for distinct integers $x,y$,
\begin{equation}\label{ten:isometry}
 v_2(\mathcal R(x)-\mathcal R(y))=v_2(x-y).
\end{equation}
Indeed, the quotient by $x-y$ is the odd integer $-1831$ plus an
even integer.  Direct substitution gives
\[
 \mathcal R(3010)\equiv0\pmod{4096},\qquad
 \mathcal R(962)\equiv2048\pmod{4096}.
\]
Equation \eqref{ten:isometry} proves that $3010$ is the unique
zero class modulo $4096$.

\begin{proposition}[The new uniform congruence range]\label{ten:range}
For $b\ge22$ and $b\not\equiv3010\pmod{4096}$,
$M_{5,b}$ is nonsingular, and
\begin{equation}\label{ten:valuation}
 v_2(\rho_{5,b})=v_2(b-3010)-5.
\end{equation}
\end{proposition}
\begin{proof}
Every $b-j$ in \eqref{ten:rho} is at least $13$.  Thus
\eqref{ten:B17} and \eqref{ten:congruence} give
\[
 \frac{P_{10}(b;2B_{2b-2},\ldots,2B_{2b-18})}{32}
          \equiv\mathcal R(b)\pmod{4096}.
\]
Outside the unique zero class, \eqref{ten:isometry} gives
$v_2(\mathcal R(b))=v_2(b-3010)<12$.
The numerator therefore has valuation $5+v_2(b-3010)$.
The denominator in \eqref{ten:rho} has valuation $10$:
$138240=2^{10}\cdot135$, and all three remaining factors are odd
at integer $b$.  This proves the nonvanishing and
\eqref{ten:valuation}.
\end{proof}

\subsection{A self-contained sufficient real-variable threshold}

We specialize the general highest-degree argument of
\cite{QuarticIncoming}; this subsection is a credited reproduction
with the explicit order-ten constants independently recomputed.
Let
\[
 J(b)=\prod_{i=0}^{4}(2n_i+1)(2n_i+2),\qquad
 E(b)=\prod_{i,\ell=0}^{4}\prod_{r=1}^{\max\{0,2i-4\ell-1\}}
                  (R_\ell+r),\qquad D(b)=10!J(b)E(b).
\]
Here $\deg E=20$ and $\deg D=30$.  Clearing the fixed factorial
and row denominators in \eqref{ten:calA} gives
\[
 \widetilde{\mathcal A}=D\mathcal A
     \in\mathbb Z[b,U_1,\ldots,U_9]^{5\times5},\qquad
 \mathscr P=\det\widetilde{\mathcal A}.
\]
Every entry has $b$-degree at most $46$, so $\deg_b\mathscr P\le230$.
If $\|f\|_1$ denotes the sum of absolute values of all coefficients,
exact expansion of the $25$ entries gives
\[
 E^*:=\max_{i,\ell}\|\widetilde{\mathcal A}_{i\ell}\|_1
 =602847821074996042728613866361551228764160.
\]
Consequently $\|\mathscr P\|_1\le K:=120(E^*)^5$.
The number $E^*$ is independently recomputed from the entry formulas
by the accompanying threshold script; it is not inferred from
numerical samples.

Put $F_0=(2b-18)!$ and
\[
 V_j=U_j\frac{F_0}{(2b-2j)!},\qquad
 \kappa=\frac{120(10!)^5}{2!4!6!8!10!}
         =14932726824960000000.
\]
The part of $\mathscr P$ homogeneous of total degree five in $U$
has the exact form
\begin{equation}\label{ten:leading}
 \mathscr P_{[5]}=-\mathscr G(b)\det(V_{i+j+1})_{0\le i,j\le4},
\end{equation}
where
\[
 \mathscr G(b)=\kappa J(b)^4E(b)^5
 \prod_{i=0}^{4}(4n_i+2)^{\underline{2i}}
 \prod_{\ell=0}^{4}\prod_{r=1}^{10+2\ell}(2b-18+r).
\]
To see the structural factorization, write
$T=\mathcal A-2F$ and use $U_j=(2b-2j)!V_j/F_0$ in
\eqref{ten:calA}.  The identity
$N_{i\ell}-(2k+1)=2b-2(5-i+\ell+k)$ cancels the factorials.
The remaining matrix factors as
\[
 T=\operatorname{diag}\!\left(\frac1{(2n_i+1)(2n_i+2)}\right)
       L\,(V_{5-i+\ell})_{0\le i,\ell\le4}
       \operatorname{diag}\!\left(\frac{R_\ell!}{F_0(2\ell+2)!}\right),
\]
where $L_{i,i-k}=d_{i,2k+1}$ is lower triangular.  Its diagonal is
$d_{i1}=-(i+1)(4n_i+2)^{\underline{2i}}$, which follows directly
from the difference of \eqref{ten:Newton} at $y=1$ and $y=0$.
Reverse the five middle rows and multiply all row factors and
$D^5$.  This gives exactly \eqref{ten:leading}.

At $U_j=2B_{2b-2j}$, Euler's Bernoulli formula yields
\begin{equation}\label{ten:actualleading}
 |\mathscr P_{[5]}|
 =\frac{\mathscr G(b)(4F_0)^5}{(2\pi)^{10b-50}}
       \det\bigl(\zeta(2b-18+2i+2j)\bigr)_{0\le i,j\le4}.
\end{equation}
The determinant is strictly positive.  Indeed, Cauchy--Binet for
the discrete moments on $m^{-2}$ gives
\[
 \det\bigl(\zeta(2b-18+2i+2j)\bigr)
 =\sum_{m_1<\cdots<m_5}(m_1\cdots m_5)^{-2b+18}
                  \prod_{r<s}(m_r^{-2}-m_s^{-2})^2.
\]
Keeping $(m_1,\ldots,m_5)=(1,2,3,4,5)$ gives the lower bound
\[
 V\,120^{-2b+18},\qquad
 V=\prod_{1\le r<s\le5}(r^{-2}-s^{-2})^2
          =\frac{2401}{36864000000000000}.
\]
For $b\ge20$ every linear factor of $\mathscr G$ is at least $b$,
and its degree is $230$, so $\mathscr G(b)\ge\kappa b^{230}$.
For $b\ge30$,
\[
 |2B_{2b-2j}|\le
 T_b:=\frac{8(2b)^{16}(2b-18)!}{(2\pi)^{2b-18}},
 \qquad 1\le j\le9,
\]
using $\zeta(2b-2j)<2$.  Also $T_b\ge1$: since $2b-18\ge42$,
the elementary comparison $26!>8^{26}$ and induction suffice.
All lower-degree terms of $\mathscr P$ have total $U$-degree at
most four.  Therefore
\[
 |\mathscr P_{[\le4]}|\le Kb^{230}T_b^4.
\]
Divide by \eqref{ten:actualleading}, cancel the powers of $2\pi$
before using $2\pi<8$, and obtain
\begin{equation}\label{ten:bound}
 \frac{|\mathscr P_{[\le4]}|}{|\mathscr P_{[5]}|}
 \le R(b):=C\frac{b^{64}960^{2b}}{(2b-18)!},\qquad
 C=\frac{K8^{26}2^{64}}{\kappa4^5V120^{18}}.
\end{equation}
Every constant is explicit and rational.  Exact integer comparison
after clearing denominators gives
\[
 R(1920)<2^{-546}<1.
\]
The verifier checks this stronger margin.  Moreover, for $b\ge1920$,
\[
 \frac{R(b+1)}{R(b)}=
 \left(1+\frac1b\right)^{64}
      \frac{960^2}{(2b-16)(2b-17)}<\frac34.
\]
For a simple uniform justification, $(1+1/b)^{64}<3$,
both denominator factors exceed $b$, and $b\ge1920=2\cdot960$.
Thus the highest-degree term cannot be cancelled for any $b\ge1920$.
Since $D(b)>0$ and the multiplier \eqref{ten:multiplier} is nonzero,
we have proved
\begin{equation}\label{ten:eventual}
             \det M_{5,b}\ne0\qquad(b\ge1920).
\end{equation}

\subsection{Completion, and an independent audit of the old exception}

For $22\le b<1920$, one cannot have $b\equiv3010\pmod{4096}$.
Proposition~\ref{ten:range} therefore proves nonvanishing throughout
this interval.  Equation \eqref{ten:eventual} supplies all
$b\ge1920$.  The remaining $6\le b\le21$ are checked directly
from the rational source matrix \eqref{ten:M}; the machine-readable
record contains every nonzero determinant.  Most of these base cases
are already contained in \cite{AlternatingIncoming}; they are recomputed here.
For example,
\[
 \det M_{5,6}=-\frac{7097067}{2867200},\qquad
 \det M_{5,8}=\frac{537653622905129}{68669669376000}.
\]
This proves Theorem~\ref{ten:main}.

The old exceptional index $b=962$ also has a short independent exact
audit.  Its numerator satisfies
\begin{equation}\label{ten:962}
 P_{10}(962;2B_{1922},2B_{1920},\ldots,2B_{1906})
       \equiv2386755584\pmod{2^{32}}.
\end{equation}
The residue has valuation $16$, giving $v_2(\rho_{5,962})=6$.
Two independent calculations of the nine Bernoulli residues agree:
the recurrence
\[
 (2m+1)(2B_{2m})=2m-1-
        \sum_{j=1}^{m-1}\binom{2m+1}{2j}(2B_{2j}),
\]
with odd division modulo $2^{32}$, and the Faulhaber power sum over
the $32768$ odd integers below $65536$, performed modulo $2^{47}$
before dividing by $2^{15}$.  Finally, evaluating the rational
five-dimensional determinant \eqref{ten:calA} directly, with
these integer residues as inputs, gives the same residue
\eqref{ten:962} as the sparse numerator.  This audit is not needed
once the uniform congruence is proved, but it verifies the exact
location where the two preceding reports failed to overlap.

\subsection{An ordinary convergent family of polygamma identities}

For every $b\ge6$ and $0\le k<b$, let
\begin{equation}\label{ten:Q}
 Q_{b,k}(x)=\sum_{t=0}^{b-1}(M_{5,b}^{-1})_{kt}
       \bigl((x+\tfrac12)^{2b+2t+1}
                    -(x-\tfrac12)^{2b+2t+1}\bigr).
\end{equation}
This form is shorter than an expansion of the even polynomial.
The identity
\[
 (x+\tfrac12)^{2j+1}-(x-\tfrac12)^{2j+1}
 =\sum_{r=0}^{j}\binom{2j+1}{2r}2^{-2(j-r)}x^{2r}
\]
shows that \eqref{ten:Q} is precisely the inverse coordinate
combination dictated by \eqref{ten:C}.

Define
\[
 \eta_{s,0}=\frac1{s-1},\qquad
 \eta_{s,j}=\frac{(s)_{2j-1}B_{2j}(1/2)}{(2j)!}\quad(j\ge1),
 \qquad c_{b,\ell}=\sum_{j=0}^{\ell}\eta_{10,j}\eta_{2b,\ell-j}.
\]
Let $S_{b,k}(x)$ be the polynomial consisting of the nonnegative
powers in the formal expansion
\[
 Q_{b,k}(x)x^{-2b-8}\sum_{\ell\ge0}c_{b,\ell}x^{-2\ell}.
\]
It is an even polynomial of degree at most $2b-10$, with rational
coefficients determined by finitely many terms.

\begin{corollary}[All order-ten reconstruction identities]
\label{ten:identity}
For every integer $b\ge6$ and $0\le k<b$,
\begin{equation}\label{ten:polygamma}
 \boxed{\displaystyle
 \sum_{n=0}^{\infty}\left\{
 Q_{b,k}(x_n)
 \frac{\psi^{(9)}(n+1)\psi^{(2b-1)}(n+1)}{9!\,(2b-1)!}
             -S_{b,k}(x_n)\right\}=\zeta(2k).}
\end{equation}
The series is ordinary and absolutely convergent.
\end{corollary}
\begin{proof}
For even $s$, $A_s(n)=\psi^{(s-1)}(n+1)/(s-1)!$.
By \eqref{ten:C}, \eqref{ten:Q}, and the exact matrix inverse,
the continued sum of $Q_{b,k}(x_n)A_{10}(n)A_{2b}(n)$ is
$\zeta(2k)$.  The centered Euler--Maclaurin expansion of the product
is exactly the formal product used to define $S_{b,k}$.
Removing all nonnegative powers leaves $O(x_n^{-2})$.
Every removed power is even and has centered continued sum
$\zeta(-2r,1/2)=0$, $r\ge0$.  Therefore the continued sum agrees
with the displayed absolutely convergent series.
\end{proof}

Here is a completely explicit member, included to make the
reconstruction mechanism concrete.  Let
\begin{align*}
 P(x)={}&58535757x^{23}+145331415x^{21}-98765059x^{19}\\
        &+92514527x^{17}-73310041x^{15}+28468113x^{13}.
\end{align*}
Then
\begin{equation}\label{ten:explicitidentity}
 \begin{split}
 \sum_{n=0}^{\infty}\Bigg\{&
    [P(n+1)-P(n)]
        [\zeta(10)-H_n^{(10)}][\zeta(12)-H_n^{(12)}]\\
   &-\frac{149591379(n+\tfrac12)^2+1835020425}{11}
                          \Bigg\}
       =-\frac{574862427}{8}.
 \end{split}
\end{equation}
It follows from the $k=0$ row of $M_{5,6}^{-1}$:
\[
 \frac4{574862427}
 (28468113,-73310041,92514527,-98765059,145331415,58535757).
\]
The numerical pair $(10,12)$ was already within the source's finite
range; the explicit identity is an illustration of the now-uniform
construction, not the new nonvanishing case.  All coefficients,
counterterms, and the exact inverse-row check are included in the
identity artifact.

\paragraph{Further questions specific to this completion.}
The unrestricted rank problem remains open in the inspected corpus
for smaller order at least twelve.  A natural next question is whether
one can prove a general overlap principle between the finite
$2$-adic exceptional classes and the constructive real-variable
threshold, without computing an increasingly large numerator.
Another is whether a determinant representation makes the lower
Bernoulli-degree terms sign-compatible, giving nonvanishing directly
and avoiding a domination threshold.  The gap-six family, where the
generic nearby-Bernoulli contraction falls outside its stable range,
still calls for a separate argument.  Finally, one may seek compact
closed forms for the inverse rows in \eqref{ten:Q}; the existence
and rational computability proved here do not by themselves make
those coefficients small.

% END sections/08_order_ten.tex

% BEGIN sections/09_cyclotomic_audit.tex
\section{A cyclotomic test case already determined by the manuscript}
\label{audit:cyclo}

The incoming alternating-tail report asks, as the first test in its
complementary-sector Question 7, for
$\Imag Z(2,1;i,-i)$, where
\[
 Z(2,1;i,-i)=\sum_{k>m\ge1}\frac{i^k(-i)^m}{k^2m}.
\]
The full two-resolvent complementary-sector generator remains open.
However, this particular coefficient is already an exact consequence of
two results in the pinned unified manuscript. Recording that implication
is an integration correction, not a newly discovered period evaluation.

\begin{proposition}[Explicit complementary coefficient]\label{audit:coefficient}
Put $\lambda_3=\Imag\Li_3((1+i)/2)$. Then
\begin{equation}\label{audit:coefficient-value}
 \boxed{\displaystyle
 \Imag Z(2,1;i,-i)
   =\lambda_3-\frac{\pi\log^2 2}{32}-\frac{\pi^3}{128}.}
\end{equation}
Equivalently,
\begin{equation}\label{audit:coefficient-integral}
 \int_0^1\frac{\log t\log(1-t)}{1+t^2}\dd t
 =\lambda_3-\frac{\pi\log^2 2}{32}-\frac{\pi^3}{128}.
\end{equation}
\end{proposition}
\begin{proof}
Use the one-direction convention
$\Li_{a,b}(z)=\sum_{k>m\ge1}z^k/(k^a m^b)$.
The inverse-color theorem in
\path{chapters/04-depth.tex}, equation
\texttt{inverse:eq:mixedb1}, gives
\[
 Z(2,1;z,z^{-1})
 =2\Li_3(z)+\Li_{2,1}(z)+\Li_{1,2}(z)-\zeta(2)\Li_1(z).
\]
The iterated-integral shuffle identity is
$\Li_1(z)\Li_2(z)=\Li_{1,2}(z)+2\Li_{2,1}(z)$.
Hence
\begin{equation}\label{audit:coefficient-reduction}
 Z(2,1;z,z^{-1})
 =2\Li_3(z)+\Li_1(z)\Li_2(z)-\Li_{2,1}(z)-\zeta(2)\Li_1(z).
\end{equation}
At $z=i$, put $L=\log2$ and $G=\beta(2)$. The standard values are
\[
 \Li_1(i)=-L/2+i\pi/4,\quad
 \Li_2(i)=-\pi^2/48+iG,\quad
 \Imag\Li_3(i)=\pi^3/32.
\]
The proved height-one specialization in
\path{chapters/04-shuffle-parity.tex} is
\[
 \Imag\Li_{2,1}(i)
 =-GL/2+\pi L^2/32-\lambda_3+3\pi^3/128.
\]
Substituting these identities into
\eqref{audit:coefficient-reduction} cancels the Catalan terms and yields
\eqref{audit:coefficient-value}. Both inherited reductions are proved in
the source, not inferred from a numerical fit \cite{ProveItPin}.

For the integral form, use the same manuscript's inverse-color integral
representation, or insert the elementary Mellin formula for $k^{-2}$
directly in the absolutely convergent double sum. It gives
\[
 Z(2,1;i,-i)
 =i\int_0^1\frac{\log t\log(1-t)}{1-it}\dd t.
\]
Taking the imaginary part proves \eqref{audit:coefficient-integral}.
The endpoint logarithms are integrable and $1-it$ does not vanish on the
integration interval, so no limiting branch choice is hidden here.
\end{proof}

The numerical value is
\[
 0.2806725061826611243960680536045272754284\ldots.
\]
This is a finite reduction into the explicitly stated classical
trilogarithm collection. It does not reduce $\lambda_3$ further and does
not prove its independence from any smaller set of constants. Nor does
it settle the full complementary generator or the distinct Gaussian
$S_6$ and revised $S_8$ conjectures.

The proposed status update is therefore precise: retain Question 7 for
the full complementary family, mark its displayed first coefficient as
already determined by the manuscript, and include
\eqref{audit:coefficient-value} as the missing cross-reference. The source
proofs of inverse-color transport and height-one reduction deserve the
credit.

% END sections/09_cyclotomic_audit.tex

% BEGIN sections/10_research_questions.tex
\section{Further research questions with concrete mathematical targets}
\label{future:sec}

The following are proposals for further work, not consequences claimed
in the preceding theorems. They prioritize exact structural statements
and finite evaluations. A question asking for a classification allows
a counterexample with a proved exceptional set as a useful answer.

\subsection{Reconstruction beyond order ten}

\begin{question}[All unequal even orders]
For every $1\le a<b$, is the first-$2b$ centered even-moment prefix
of the $(2a,2b)$ tail product nonsingular in its formal source
coordinates? The present article completes $a=5$; the project now
has all $a\le5$. A decisive next result would be a structural
noncancellation theorem for the fixed-size Bernoulli determinant,
or a proved description of any exceptional pairs.
\end{question}

\begin{question}[A uniform overlap principle]
Can one relate the precision of a doubled-Bernoulli congruence to the
constructive real-variable threshold for arbitrary fixed $a$?
At $a=5$ the sole exceptional class has least positive representative
$3010$, exceeding the threshold $1920$. Higher precision alone does
not prove nonvanishing: a compatible $2$-adic root may persist. The
target is a theorem ensuring coverage by two complementary arguments,
not a table of increasingly long nonzero prefixes.
\end{question}

\begin{question}[The full gap-six family]
Give a bordered-Hankel reduction that retains the moving $B_0$ and
$B_1$ endpoint terms for $b=a+3$. The stable fixed-order contraction
requires $b\ge2a$ and therefore does not cover this family uniformly.
A formula with all boundary terms retained could settle the global
gap-six conjecture without extrapolating the finitely checked pairs.
\end{question}

\begin{question}[Compact inverse rows]
Find a product, determinant of uniformly smaller size, or short
recurrence for the polynomials $Q_{b,k}$ of
Corollary~\ref{ten:identity}. Their exact rational computability is
proved, but the displayed $(10,12)$ example shows that naive expanded
coefficients grow rapidly. A compact expression should also identify
the subtraction polynomial directly and preserve ordinary convergence.
\end{question}

\subsection{Beyond the bilinear Gamma normalization}

\begin{question}[Every bilinear contact hyperplane]
For the two independent orders in Theorem~\ref{bil:master}, compute
the constant terms of $C_N$ and $(-1)^ND_N$ on the further resonant
hyperplanes $p+q=-2m$. The first symmetric contact is $1/2$ at
$p+q=0$. Seek closed bivariate contact polynomials, every independent
$p,q$ jet, and the exact change under translating the cutoff.
The diagonal answer in Section~\ref{cut:sec} is a mandatory check;
it does not determine all transverse derivatives.
\end{question}

\begin{question}[Independent outer weights and ordinary Hurwitz tails]
Construct a finite-cutoff identity for
$\zeta(p,n+a)\zeta(q,n+b)(n+c)^{-w}$ with independent
$p,q,w,a,b,c$, retaining an explicit remainder or a precisely defined
finite collection of additional functions. Determine its moving
polar hyperplanes and residues before differentiating on them.
The Gamma-weighted alternating identities do not settle this ordinary
Hurwitz problem; a valid result must state the size and analytic
meaning of any newly required coordinate collection.
\end{question}

\begin{question}[More colors and higher products]
For tails with a primitive $d$th-root twist, can a product of Gamma
quotients convert the $d$-step recurrence into a finite bilinear
endpoint law? For three tail factors, identify the smallest adjacent
product ansatz whose linear and quadratic residual terms can all
cancel. A useful target is a divisor-orbit criterion extending
Theorem~\ref{rat:classification}, together with one ordinary convergent
three-factor evaluation that genuinely lies outside its bilinear case.
\end{question}

\begin{question}[The remaining logarithmic Gamma generator]
The logarithmic square sum in Corollary~\ref{cut:Gamma} is explicit,
but the earlier twisted logarithmic-Gamma correction generator has
an independent outer resolvent. Determine whether it belongs to a
finite stated collection of classical or multiple-Gamma functions,
or prove a functional obstruction in a declared difference-equation
class. A failure of numerical reduction is not such an obstruction.
\end{question}

\subsection{Harmonic pencils, spectral coefficients, and polylogarithms}

\begin{question}[Complex weights and sharp cutoff exceptions]
Does generic degree--rank equality extend to every complex rational
weight? If not, classify the smallest counterexamples using the
kernel-compression identity \eqref{deg:highest}. For real weights,
determine whether the bound $r(r-\rho)$ on late deficient cutoffs is
sharp in the rational-moment class, and how it improves for special
polynomial kernel spaces. The real proof uses definite increments;
it cannot be transferred to a complex bilinear Gram form by wording
alone.
\end{question}

\begin{question}[Higher-rank pencils and their realizations]
Classify rational residue data whose entire harmonic pencil has
generic rank two, then rank $\rho$. The family in
Theorem~\ref{hier:rank} realizes every rank, but does not classify all
pencils: the maximal-rank kernel need not be a common kernel.
Seek exact finite residue conditions and identify which such pencils
come from positive rational weights. This would complement the
existing rank-one classification.
\end{question}

\begin{question}[All Laurent coefficients in compact form]
For the hierarchy $Z_{r,\rho}(s,a)$, find short formulas or recurrences
for the lower normalizations $G_{L,h}^{(m)}(a)$ at $m\ge2$, and for
the complete principal and finite parts at every pole. The leading
coefficient at $m=1$ is a Gamma product. At higher orders, confluent
deformations of the mixed factorial moments may reduce determinant
growth. Any claimed formula should reproduce both finite parts in
\eqref{hier:r2-s1}--\eqref{hier:r2-s2} and distinguish them from
hard-cutoff constants.
\end{question}

\begin{question}[Complementary cyclotomic generators and exact depth]
Construct the full complementary two-resolvent generator requested
in the alternating report, retaining the evaluated coefficient
\eqref{audit:coefficient-value} as a boundary check. In parallel,
determine minimal depth for concrete harmonic-determinant generating
functions inside a specified shuffle, differential, or symbol
calculus. Formal harmonic degree predicts a useful representation
bound; it does not by itself prove minimal polylogarithmic depth.
The Gaussian $S_6$ and revised $S_8$ assertions remain separate
conjectures requiring exact reductions or counterexamples.
\end{question}

\subsection{A practical order of attack}

The next finite target is smaller order twelve: derive its exact
fixed-size numerator, determine a sufficiently precise Bernoulli
congruence, and compare the exceptional classes with the real-variable
threshold. The next analytic target is the first off-diagonal contact
hyperplane $p+q=-2$, where the new bilinear endpoint contains enough
independent parameters to reveal which diagonal parity cancellations
survive. The next determinant target is a rank-two residue
classification, tested against both the positive shorter-stencil
family and a pencil without a common kernel.

For formal verification, the finite bilinear recurrence, rational
root-orbit criterion, and symmetric-pencil compression lemma are
small natural starting points. The order-ten congruence can then be
checked in a polynomial ring modulo $2^{17}$ from the finite Newton
matrix. Such formalization would strengthen the trust chain without
altering the claimed mathematics or replacing a proof by decimal
agreement.

% END sections/10_research_questions.tex

% BEGIN sections/11_verification.tex
\section{Verification, source corrections, and integration}
\label{verification:sec}

\subsection{Proofs and exact computational certificates}

The bilinear, rational-certificate, cutoff-contact, harmonic-degree,
and spectral-hierarchy theorems have proofs in the article. Their
finite formulas also admit independent symbolic and numerical
replays. The order-ten completion additionally uses an exact finite
computer-assisted certificate: a polynomial defined by a $5\times5$
Newton/factorial determinant is reconstructed coefficientwise, then
substituted into an explicit congruence in $\Z[b]/(2^{17})$.
All of the inputs, finite algorithms, and integer comparisons are
included. The proof does not depend on floating-point nonzero tests.

The complete polynomial reconstruction sums all $120$ signed
permutation products and matches $10\,403$ expanded coefficients,
stored in $439$ groups. The old exceptional integer $b=962$ is
checked separately by two Bernoulli computations and by direct
evaluation of the $5\times5$ determinant. The sufficient threshold
$b=1920$ is regenerated from the source entry formulas and compared
using exact integers. Finite source matrices handle $6\le b\le21$.

\subsection{Reproducible diagnostics}

The included programs use Python, SymPy, and mpmath; the observed
versions were Python 3.12, SymPy 1.14.0, and mpmath 1.3.0. They need
no network access. The package README supplies commands and file
locations. The separate result records preserve the precision and
test method rather than presenting all checks as equivalent evidence.

\begingroup\small
\begin{longtable}{@{}>{\raggedright\arraybackslash}p{.24\textwidth}>{\raggedright\arraybackslash}p{.42\textwidth}>{\raggedright\arraybackslash}p{.28\textwidth}@{}}
\toprule
Component & Exact checks or independent route & Role \\
\midrule\endfirsthead
\toprule Component & Exact checks or independent route & Role \\
\midrule\endhead
Bilinear tails & Five symbolic cancellations; 28 numerical checks,
including independent Laplace integrals for Stieltjes jets. &
Recurrence-sign, normalization, domain, and example diagnostics. \\
\addlinespace
Cutoff jets and cyclotomic audit & 74 exact Euler-polynomial and
translated-coefficient checks through order twelve; 12 numerical
checks, including independent quadrature. & Contact signs, sharp
threshold, differentiated finite sums, and the inherited trilogarithm
coefficient. \\
\addlinespace
Degree and spectral hierarchy & 45 exact checks, including all
ranks for $r=2,\ldots,6$; 19 numerical checks through independent
Mellin moments and Laurent limits. & Non-common-kernel example,
signed cutoff exception, pole normalization, Gamma products, and
Stieltjes finite parts. \\
\addlinespace
Polylogarithm generator & 122 exact cutoff and Taylor-coefficient
checks; independent endpoint evaluation from direct moments and
harmonic tail expansions. & Generator, anchored primitive, and
ordinary-zeta endpoint; 65-digit numerical diagnostics. \\
\addlinespace
Order-ten completion & Full coefficient reconstruction, uniform
congruence, exact base matrices and threshold; two independent
Bernoulli residue calculations. & Exact finite certificates supporting
the uniform nonvanishing theorem. \\
\addlinespace
Explicit $(10,12)$ identity & Exact inverse-row and counterterm
arithmetic; independent numerical cutoff/tail diagnostics. & An
ordinary convergent illustration of the uniform family. \\
\bottomrule
\end{longtable}
\endgroup

The 45-digit bilinear checks have residuals below $3.1\cdot10^{-45}$.
Independent Mellin comparisons for the determinant formulas agree
at roughly $10^{-79}$ relative error at 80-digit precision;
the finite-step Laurent extrapolations have residuals of order
$10^{-24}$. The polylogarithm endpoint checks have errors below
$2\cdot10^{-41}$ and $4\cdot10^{-48}$ with two truncation orders.
The displayed $(10,12)$ identity has numerical residuals
below $4\cdot10^{-68}$. These are numerical diagnostics, not certified
interval bounds. They neither prove global identities nor establish
independence of constants. The written arguments and exact algebraic
certificates have that separate role.

Independent mathematical audits checked the bilinear endpoint
normalization, the rational factor $cS(0)$, translated cutoff jets,
the pencil compression argument, the finite-difference residue
blocks, and the Gamma-product normalization. No unresolved
discrepancy was found in those checks. This records the verification
performed; it does not describe the article as externally refereed
or formally verified in a proof assistant.

\subsection{Corrections and status updates proposed for ProveIt}

The concrete correction is to the first coefficient in the incoming
complementary-sector question. Add the formula in
Proposition~\ref{audit:coefficient} and cross-reference the existing
inverse-color and height-one manuscript reductions. Retain the full
generator question. The correction concerns an overlooked implication,
not an incorrect proved identity.

The order-ten progression should now be marked proved by
Theorem~\ref{ten:main}; its prior congruence remains valid on its
original range. The new result does not invalidate that earlier proof.
The degree--rank question should be split into the proved real-rational
generic theorem, the explicit finite-cutoff exception test, and the
remaining complex/depth questions. The ordinary-square continuation
question is resolved with its cutoff coordinate included in the
statement.

Two presentation precautions are mathematically necessary. First,
``odd negative orders have no contact'' is true for $X=N+a+1$;
it is false for an arbitrary translated cutoff. Second, all-coordinate
formal rank is not arithmetic independence of zeta values. These
qualifications should travel with the imported statements.

No new error in the already proved core manuscript is asserted by
this audit, and this work is not an exhaustive audit of the whole
repository. The Gaussian $S_6$ and revised $S_8$ statements are not
proved here. The unrestricted even-order prefix theorem, complete
gap-six family, ordinary independent-order outer-weight problem,
and full complementary generator remain explicitly open.

\subsection{Integration boundary}

The package supplies a modular \path{article.tex}, section sources,
a flattened standalone TeX file, the rendered PDF, exact verification
programs and records, and source/status notes. Section labels have
disjoint prefixes so individual results can be imported. The supplied
source pin and question map identify what was inspected. The work
does not modify the repository, supersede all earlier reports, or
bundle their generated build files.

% END sections/11_verification.tex

% BEGIN references.tex
\begin{thebibliography}{99}
\raggedright

\bibitem{ProveItPin}
V. Reshetnikov, ProveIt repository, unified polylogarithm manuscript.
Inspected revision
\texttt{da2792fc9cc8ef19c6a314be6c65e1437fbdc66b},
11 October 2026.
\url{https://github.com/VladimirReshetnikov/ProveIt}.
Relevant manuscript files are
\path{Analysis/Polylogarithms/docs/manuscript/chapters/04-depth.tex}
and \path{chapters/04-shuffle-parity.tex} relative to that manuscript.

\bibitem{AlternatingIncoming}
ProveIt research corpus,
\emph{Alternating Tails and Higher Spectral Defects},
incoming package
\path{ProveIt_Alternating_Tails_Spectral_Defects_2026-10-11.zip},
at the revision in \cite{ProveItPin}.
Sections on polynomial adjacent-tail certificates and order-ten
congruences; further-research Questions 1--7.

\bibitem{FiniteDifferenceIncoming}
ProveIt research corpus, \emph{Finite Difference Determinants},
incoming package
\path{ProveIt_Finite_Difference_Determinants_2026-10-11.zip},
at the revision in \cite{ProveItPin}.
Relevant sections:
\path{sections/02-spectral.tex},
\path{sections/06-classification.tex},
\path{sections/09-questions.tex}, and
\path{sections/10-residue-product.tex}.

\bibitem{QuarticIncoming}
ProveIt research corpus,
\emph{Harmonic-Tail Reconstruction, Quartic Trigamma Moments,
and Spectral Stieltjes Calculus},
incoming package
\path{ProveIt_Quartic_Spectral_Reconstruction_2026-10-11.zip},
at the revision in \cite{ProveItPin}.
Sections 2--3 supply the inherited order-eight completion and the
general constructive eventual-rank argument.

\bibitem{GeneralHankel}
ProveIt research corpus, \emph{Bernoulli--Tornheim--Digamma continuation},
``A fixed-size determinant for every pair of even orders,''
\path{sections/03b_general_hankel.tex}, in
\path{Analysis/Polylogarithms/docs/reports/stieltjes-harmonic-resonance/continuations/bernoulli-tornheim-digamma/},
at the revision in \cite{ProveItPin}.

\bibitem{HigherQuadraticIncoming}
ProveIt research corpus,
incoming package
\path{ProveIt_Higher_Quadratic_Jets_2026-10-11.zip},
at the revision in \cite{ProveItPin}.
Inspected as prior context; its higher-order jet constructions are
not claimed as new in the present article.

\bibitem{HigherSpectralIncoming}
ProveIt research corpus,
incoming package
\path{ProveIt_Higher_Spectral_Defects_2026-10-11.zip},
at the revision in \cite{ProveItPin}.
Inspected as prior context; the present shorter-stencil family is
distinguished from its general spectral-defect questions.

\bibitem{bil:HuKim}
S. Hu and M.-S. Kim,
\emph{On the Stieltjes constants and gamma functions with respect to
alternating Hurwitz zeta functions},
J. Math. Anal. Appl. \textbf{509} (2022), no. 1, article 125930.
\url{https://arxiv.org/abs/2106.14674}.

\bibitem{bil:WangHuKim}
W. Wang, S. Hu, and M.-S. Kim,
\emph{On gamma functions with respect to the alternating Hurwitz zeta
functions}, arXiv:2405.02854v4 (2024).
\url{https://arxiv.org/abs/2405.02854v4}.

\bibitem{HuKimAsymptotic}
S. Hu and M.-S. Kim,
\emph{Asymptotic expansions for the alternating Hurwitz zeta function
and its derivatives}, arXiv:2103.15528v6 (2023).
\url{https://arxiv.org/abs/2103.15528}.
The underlying uniform expansions, Euler specializations, and
spectral derivative calculus are classical input here.

\bibitem{bil:Hoffman}
M. E. Hoffman,
\emph{Sums of products of Riemann zeta tails},
Mediterr. J. Math. \textbf{13} (2016), 2771--2781.
\url{https://doi.org/10.1007/s00009-015-0653-9};
\url{https://arxiv.org/abs/1602.03399}.

\bibitem{rat:Abramov}
S. A. Abramov,
\emph{The rational component of the solution of a first-order linear
recurrence relation with rational right-hand side},
U.S.S.R. Comput. Math. Math. Phys. \textbf{15} (1975), no. 4,
216--221.
\url{https://doi.org/10.1016/0041-5553(75)90181-0}.

\bibitem{DLMFZeta}
NIST Digital Library of Mathematical Functions, Section 25.11,
Hurwitz zeta function.
\url{https://dlmf.nist.gov/25.11}.

\bibitem{bil:DLMF}
NIST Digital Library of Mathematical Functions, Sections 5.5 and 5.11,
Gamma functional relations and asymptotic expansions.
\url{https://dlmf.nist.gov/5.5};
\url{https://dlmf.nist.gov/5.11}.

% Add this item inside the article's thebibliography environment.
\bibitem{hier:BB2005}
Jonathan M. Borwein and David M. Bradley,
\emph{On Two Fundamental Identities for Euler Sums},
preprint, January 2005, especially Eq.~(1.3).
\url{https://carmamaths.org/resources/jon/Preprints/Papers/ArXiv/Zeta21/z21.pdf}.

\end{thebibliography}

% END references.tex

\end{document}
